%%%%%%%%%%%%%%%%%%%%%%%%%%%%%%%%%%%%%%%%%
%Ayana's thesis template
%Version 1.0 (03/06/2015)
%Author: Ayana B. Martins (ayanamartins(at)gmail.com)

%Based on:
% The Legrand Orange Book (original author: Mathias Legrand)
% Version 1.3 (21/8/13)
% Downloaded from LaTeXTemplates: http://www.latextemplates.com/template/the-legrand-orange-book

%%%%%%%%%%%%%%%%%%%%%%%%%%%%%%%%%%%%%%%%%%
% License:
% CC BY-NC-SA 3.0 (http://creativecommons.org/licenses/by-nc-sa/3.0/)
%
% Compilation
% This template uses biblatex for its bibliography and makeindex for its index.% 
% When you first open the template, compile it from the command line with the 
% commands below to make sure your LaTeX distribution is configured correctly:
%
% 1) pdflatex thesis
% 2) makeindex thesis.idx -s StyleInd.ist
% 3) for i in {1..N}; do bibtex thesis${i}-blx ; done
% where N is the number of refsections (chapters, 3 in 
%this case)
% 4) pdflatex thesis (x 2)
%
% After this, when you wish to update the bibliography/index use the appropriate
% command above and make sure to compile with pdflatex several times 
% afterwards to propagate your changes to the document.
%
% This template also uses a number of packages which may need to be
% updated to the newest versions for the template to compile. It is strongly
% recommended you update your LaTeX distribution if you have any
% compilation errors.
%
% Important note:
% Chapter heading images should have a 2:1 width:height ratio,
% e.g. 920px width and 460px height.
%
%%%%%%%%%%%%%%%%%%%%%%%%%%%%%%%%%%%%%%%%%

%----------------------------------------------------------------------------------------
%	PACKAGES AND OTHER DOCUMENT CONFIGURATIONS
%----------------------------------------------------------------------------------------

\documentclass[11pt, % Default font size
               %fleqn, %left-justified equations
               oneside] %one-sided or two-sided documents
               {book}
\usepackage[top=3cm,bottom=3cm,left=3.2cm,right=3.2cm,headsep=10pt,a4paper]{geometry} % Page margins
\usepackage{xcolor} % Required for specifying colors by name
\definecolor{color}{RGB}{128,0,0} % Define the color used for highlighting throughout the book
\definecolor{blue}{RGB}{69,74,159} %Redifines the default for the color blue
\usepackage{lastpage}

%----------------------------------------------------------------------------------------
% Font Settings
\usepackage{avant} % Use the Avantgarde font for headings
%\usepackage{times} % Use the Times font for headings
\usepackage{mathptmx} % Use the Adobe Times Roman as the default text font together with math symbols from the Sym­bol, Chancery and Com­puter Modern fonts
\usepackage{microtype} % Slightly tweak font spacing for aesthetics
\usepackage[utf8]{inputenc} % Required for including letters with accents
\usepackage[T1]{fontenc} % Use 8-bit encoding that has 256 glyphs
%\fontfamily{arial}\selectfont 
\usepackage{csquotes}
\usepackage{endnotes}
\usepackage{setspace}
\onehalfspacing


%-----------------------------------------
%Hyperrref is used in the template only
%to link some of the packages that are called
\usepackage[
  implicit = false,
  bookmarks=false,
  pdfpagelabels=false,
  hyperfootnotes=false,
  hyperindex=false,
  colorlinks,
  urlcolor=blue
]{hyperref}

\makeatletter
\let\saved@hyper@linkurl\hyper@linkurl
%\let\saved@hyper@linkfile\hyper@linkfile
\let\saved@hyper@link@\hyper@link@
\AtBeginDocument{%
  % Since the whole document is affected, only the \begin part of
  % environment `NoHyper' is needed.
  \NoHyper
  \let\hyper@linkurl\saved@hyper@linkurl % needed by \url
  %\let\hyper@linkfile\saved@hyper@linkfile % needed by \href{<file>}
  \let\hyper@link@\saved@hyper@link@ % needed by \href{<url>}
}
\makeatother

%----------------------------------------------------------------------------------------
% Bibliography

\usepackage[style=custom-authoryear,
            %sorting=nyt,
            %sortcites=true,
            autopunct=true,
            %babel=hyphen,
            hyperref=true,
            %abbreviate=false,
            firstinits=true,
            backref=true,
            refsection=chapter,
            url=false, doi=false, isbn=false,
            backend=bibtex]
{biblatex}
% back reference text preceding the page number ("see p.")
\DefineBibliographyStrings{english}{%
    backrefpage  = {see p.}, % for single page number
    backrefpages = {see pp.} % for multiple page numbers
}

% the followings activate 'custom-english-ordinal-sscript.lbx'
% in order to print ordinal 'edition' suffixes as superscripts,
% and adjusts (reduces) spacing between suffix and following "ed."
\DeclareLanguageMapping{english}{custom-english-ordinal-sscript}
\DeclareFieldFormat{edition}%
                   {\ifinteger{#1}%
                    {\mkbibordedition{#1}\addthinspace{}ed.}%
                    {#1\isdot}}

% removes period at the very end of bibliographic record
\renewcommand{\finentrypunct}{}

% removes period after DOI and suppresses capitalization
% of the word following DOI ("See p. xx" -> "see p. xx")
\renewcommand{\newunitpunct}{\addspace\midsentence}

\DeclareFieldFormat{journaltitle}{\mkbibemph{#1},} % italic journal title with comma
\DeclareFieldFormat[inbook,thesis]{title}{\mkbibemph{#1}\addperiod} % italic title with period
\DeclareFieldFormat[article]{title}{#1} % title of journal article is printed as normal text
\DeclareFieldFormat[article]{volume}{\textbf{#1}\addcolon\space} % makes volume of journal bold and adds colon
\DeclareFieldFormat{pages}{#1} % removes pagination (p./pp.) before page numbers

%%%%%%%%%
% the command \sjcitep defined below prints footnote citation above punctuation
\newlength{\spc} % declare a variable to save spacing value
\newcommand{\sjcitep}[2][]{% new command with two arguments: optional (#1) and mandatory (#2)
        \settowidth{\spc}{#1}% set value of \spc variable to the width of #1 argument
        \addtolength{\spc}{-1.8\spc}% subtract from \spc about two (1.8) of its values making its magnitude negative
        #1% print the optional argument
        \hspace*{\spc}% print an additional negative spacing stored in \spc after #1
        \supershortnotecite{#2}}% print (cite) the mandatory argument
%%%%%%%%% 
\addbibresource{example.bib} % BibTeX bibliography file
\defbibheading{bibempty}{}


%Another option with natbib which require a different method to do a bibliography section per chapter
%\usepackage[square,numbers,sectionbib]{natbib}
%\usepackage{chapterbib}
%\usepackage{authblk}
%\bibliographystyle{evolution}

%----------------------------------------------------------------------------------------
% Index
\usepackage{calc} % For simpler calculation - used for spacing the index letter headings correctly
\usepackage{makeidx} % Required to make an index
\makeindex % Tells LaTeX to create the files required for indexing

%----------------------------------------------------------------------------------------
%----------------------------------------------------------------------------------------
%	VARIOUS REQUIRED PACKAGES
%----------------------------------------------------------------------------------------

\usepackage{titlesec} % Allows customization of titles
\usepackage{graphicx} % Required for including pictures
\graphicspath{{Images/}} % Specifies the directory where pictures are stored
\usepackage{lipsum} % Inserts dummy text
\usepackage{tikz} % Required for drawing custom shapes
\usepackage[brazil]{babel} % Specify language/hyphenation
\usepackage{enumitem} % Customize lists
\setlist{nolistsep} % Reduce spacing between bullet points and numbered lists
\usepackage{booktabs} % Required for nicer horizontal rules in tables
\usepackage{eso-pic} % Required for specifying an image background in the title page
\usepackage[rightcaption]{sidecap} % Required for placing caption on the side of float 

%----------------------------------------------------------------------------------------
%	MAIN TABLE OF CONTENTS
%----------------------------------------------------------------------------------------

\usepackage{titletoc} % Required for manipulating the table of contents

\addtocontents{toc}{\protect\thispagestyle{empty}}
\contentsmargin{0cm} % Removes the default margin
% Chapter text styling
\titlecontents{chapter}[1.25cm] % Indentation
{\addvspace{15pt}\large\sffamily\bfseries} % Spacing and font options for chapters
{\color{color!60}\contentslabel[\Large\thecontentslabel]{1.25cm}\color{color}} % Chapter number
{}  
{\color{color!60}\normalsize\sffamily\bfseries\;\titlerule*[.5pc]{.}\;\thecontentspage} % Page number
% Section text styling
\titlecontents{section}[1.25cm] % Indentation
{\addvspace{5pt}\sffamily\bfseries} % Spacing and font options for sections
{\contentslabel[\thecontentslabel]{1.25cm}} % Section number
{}
{\sffamily\hfill\color{black}\thecontentspage} % Page number
[]
% Subsection text styling
\titlecontents{subsection}[1.25cm] % Indentation
{\addvspace{1pt}\sffamily\small} % Spacing and font options for subsections
{\contentslabel[\thecontentslabel]{1.25cm}} % Subsection number
{}
{\sffamily\;\titlerule*[.5pc]{.}\;\thecontentspage} % Page number
[] 

%----------------------------------------------------------------------------------------
%	MINI TABLE OF CONTENTS IN CHAPTER HEADS
%----------------------------------------------------------------------------------------

% Section text styling
\titlecontents{lsection}[0em] % Indendating
{\footnotesize\sffamily} % Font settings
{}
{}
{}

% Subsection text styling
\titlecontents{lsubsection}[.5em] % Indentation
{\normalfont\footnotesize\sffamily} % Font settings
{}
{}
{}
 
%----------------------------------------------------------------------------------------
%	PAGE HEADERS
%----------------------------------------------------------------------------------------

\usepackage{fancyhdr} % Required for header and footer configuration

\pagestyle{fancy}
\pagenumbering{roman} 
\renewcommand{\chaptermark}[1]{\markboth{\sffamily\normalsize\bfseries #1}{}} % Chapter text font settings
\renewcommand{\sectionmark}[1]{\markright{\sffamily\normalsize\thesection\hspace{5pt}#1}{}} % Section text font settings


%-----------------------------------------------------------------------------------------
%Two-sided printing
%\fancyhf{} \fancyhead[LE,RO]{\sffamily\normalsize\thepage} % Font setting for the page number in the header
%\fancyhead[LO]{\rightmark} % Print the nearest section name on the left side of odd pages
%\fancyhead[RE]{\leftmark} % Print the current chapter name on the right side of even pages
%------------------------------------------------------------------------------------------

%-----------------------------------------------------------------------------------------
%One-sided printing
\fancyhf{} \fancyhead[RE,RO]{\sffamily\normalsize\thepage} % Font setting and placement for the page number in the header
\fancyhead[LO]{\rightmark} % Print the nearest section name on the left side of odd pages
\fancyhead[LE]{\leftmark} % Print the current chapter name on the right side of even pages
%------------------------------------------------------------------------------------------

\renewcommand{\headrulewidth}{0.5pt} % Width of the rule under the header
\addtolength{\headheight}{2.5pt} % Increase the spacing around the header slightly
\renewcommand{\footrulewidth}{0pt} % Removes the rule in the footer
\fancypagestyle{plain}{\fancyfoot[CO,CE]{\thepage}
                      \renewcommand{\headrulewidth}{0pt}} % Style for when a plain pagestyle is specified


\makeatletter
% Removes the header from odd empty pages at the end of chapters
%\renewcommand{\cleardoublepage}{
%\clearpage\ifodd\c@page\else
%\hbox{}
%\vspace*{\fill}
%\thispagestyle{empty}
%\newpage
%\fi}

%----------------------------------------------------------------------------------------
%	SET UP POSITION OF TEXT IN TITLE PAGE
%----------------------------------------------------------------------------------------
\def\parsecomma#1,#2\endparsecomma{\def\page@x{#1}\def\page@y{#2}}
\tikzdeclarecoordinatesystem{page}{
    \parsecomma#1\endparsecomma
    \pgfpointanchor{current page}{north east}
    % Save the upper right corner
    \pgf@xc=\pgf@x%
    \pgf@yc=\pgf@y%
    % save the lower left corner
    \pgfpointanchor{current page}{south west}
    \pgf@xb=\pgf@x%
    \pgf@yb=\pgf@y%
    % Transform to the correct placement
    \pgfmathparse{(\pgf@xc-\pgf@xb)/2.*\page@x+(\pgf@xc+\pgf@xb)/2.}
    \expandafter\pgf@x\expandafter=\pgfmathresult pt
    \pgfmathparse{(\pgf@yc-\pgf@yb)/2.*\page@y+(\pgf@yc+\pgf@yb)/2.}
    \expandafter\pgf@y\expandafter=\pgfmathresult pt
}

%----------------------------------------------------------------------------------------
%	MODIFIED ENDNOTE AND FOOTNOTES COMMANDS
%----------------------------------------------------------------------------------------
\renewcommand{\theendnote}{\alph{endnote}}

\newcommand\footnoteref[1]{\protected@xdef\@thefnmark{\ref{#1}}\@footnotemark}

%Use the following command to add a endnote to a footnote without showing the number.

\newcommand\numberlessnote[1]{%
  \begingroup
  \renewcommand\thefootnote{}\footnote{#1}%
  \addtocounter{footnote}{-1}%
  \endgroup
}

%\newcommand\insidenote[1]{%
%  \begingroup
%  \addtocounter{endnote}{+1}%
%  \renewcommand\theendnotes{}\endnotetext{#1}%
%  \endgroup
%}

%Use the following command if you are using insidenotes in our text and one to keep which
%note number unique
%Example:
% Some text \footnote{See page \ref{endnotes} \insidenote{Actual note}}
% The label endnote should then be placed in the page with the notes.

%\newcommand\onlyfootnote[1]{%
%  \begingroup
%  \addtocounter{endnote}{+1}%
%  \renewcommand\theendnotes{}\footnote{#1}%
%  \endgroup
%}

%\renewcommand*{\thefootnote}{\fnsymbol{footnote}}

%----------------------------------------------------------------------------------------
%	MODIFIED TABLES
%----------------------------------------------------------------------------------------

\usepackage{tabularx}
\newcolumntype{C}[1]{>{\hsize=#1\hsize\centering\arraybackslash}X}%
\newcolumntype{L}[1]{>{\hsize=#1\hsize\raggedright\arraybackslash}X}%


\def\hlinewd#1{%
\noalign{\ifnum0=`}\fi\hrule \@height #1 %
\futurelet\reserved@a\@xhline}


\newcommand{\abstractseparator}{
\textcolor{color}{\rule[0.5ex]{1\columnwidth}{2pt}}
}

%----------------------------------------------------------------------------------------
%	THEOREM STYLES
%----------------------------------------------------------------------------------------


\usepackage{amsmath,amsfonts,amssymb,amsthm} % For including math equations, theorems, symbols, etc

\newcommand{\intoo}[2]{\mathopen{]}#1\,;#2\mathclose{[}}
\newcommand{\ud}{\mathop{\mathrm{{}d}}\mathopen{}}
\newcommand{\intff}[2]{\mathopen{[}#1\,;#2\mathclose{]}}
\newtheorem{notation}{Notation}[chapter]

%%%%%%%%%%%%%%%%%%%%%%%%%%%%%%%%%%%%%%%%%%%%%%%%%%%%%%%%%%%%%%%%%%%%%%%%%%%
%%%%%%%%%%%%%%%%%%%% dedicated to boxed/framed environements %%%%%%%%%%%%%%
%%%%%%%%%%%%%%%%%%%%%%%%%%%%%%%%%%%%%%%%%%%%%%%%%%%%%%%%%%%%%%%%%%%%%%%%%%%

\newtheoremstyle{ocrenumbox}% % Theorem style name
{0pt}% Space above
{0pt}% Space below
{\normalfont}% % Body font
{}% Indent amount
{\small\bf\sffamily\color{color}}% % Theorem head font
{\;}% Punctuation after theorem head
{0.25em}% Space after theorem head
{\small\sffamily\color{color}\thmname{#1}\nobreakspace\thmnumber{\@ifnotempty{#1}{}\@upn{#2}}% Theorem text (e.g. Theorem 2.1)
\thmnote{\nobreakspace\the\thm@notefont\sffamily\bfseries\color{black}---\nobreakspace#3.}} % Optional theorem note
\renewcommand{\qedsymbol}{$\blacksquare$}% Optional qed square

\newtheoremstyle{blacknumex}% Theorem style name
{5pt}% Space above
{5pt}% Space below
{\normalfont}% Body font
{} % Indent amount
{\small\bf\sffamily}% Theorem head font
{\;}% Punctuation after theorem head
{0.25em}% Space after theorem head
{\small\sffamily{\tiny\ensuremath{\blacksquare}}\nobreakspace\thmname{#1}\nobreakspace\thmnumber{\@ifnotempty{#1}{}\@upn{#2}}% Theorem text (e.g. Theorem 2.1)
\thmnote{\nobreakspace\the\thm@notefont\sffamily\bfseries---\nobreakspace#3.}}% Optional theorem note

\newtheoremstyle{blacknumbox} % Theorem style name
{0pt}% Space above
{0pt}% Space below
{\normalfont}% Body font
{}% Indent amount
{\small\bf\sffamily}% Theorem head font
{\;}% Punctuation after theorem head
{0.25em}% Space after theorem head
{\small\sffamily\thmname{#1}\nobreakspace\thmnumber{\@ifnotempty{#1}{}\@upn{#2}}% Theorem text (e.g. Theorem 2.1)
\thmnote{\nobreakspace\the\thm@notefont\sffamily\bfseries---\nobreakspace#3.}}% Optional theorem note


\newtheoremstyle{coloredbox}% Box style name
{0pt}% Space above
{0pt}% Space below
{\normalfont}% % Body font
{}% Indent amount
{\small\bf\sffamily\color{color}}% % Theorem head font
{\;}% Punctuation after theorem head
{\newline}% Space after theorem head
{\small\sffamily\color{color}\thmname{#1}\nobreakspace\thmnumber{\@ifnotempty{#1}{}\@upn{#2}}% Theorem text (e.g. Theorem 2.1)
\thmnote{\nobreakspace\the\thm@notefont\sffamily\bfseries\color{black}---\nobreakspace#3.}} % Optional theorem note
\renewcommand{\qedsymbol}{$\blacksquare$}% Optional qed square

%%%%%%%%%%%%%%%%%%%%%%%%%%%%%%%%%%%%%%%%%%%%%%%%%%%%%%%%%%%%%%%%%%%%%%%%%%%
%%%%%%%%%%%%% dedicated to non-boxed/non-framed environements %%%%%%%%%%%%%
%%%%%%%%%%%%%%%%%%%%%%%%%%%%%%%%%%%%%%%%%%%%%%%%%%%%%%%%%%%%%%%%%%%%%%%%%%%
\newtheoremstyle{ocrenum}% % Theorem style name
{5pt}% Space above
{5pt}% Space below
{\normalfont}% % Body font
{}% Indent amount
{\small\bf\sffamily\color{color}}% % Theorem head font
{\;}% Punctuation after theorem head
{0.25em}% Space after theorem head
{\small\sffamily\color{color}\thmname{#1}\nobreakspace\thmnumber{\@ifnotempty{#1}{}\@upn{#2}}% Theorem text (e.g. Theorem 2.1)
\thmnote{\nobreakspace\the\thm@notefont\sffamily\bfseries\color{black}---\nobreakspace#3.}} % Optional theorem note
\renewcommand{\qedsymbol}{$\blacksquare$}% Optional qed square
\makeatother

% Defines the theorem text style for each type of theorem to one of the three styles above
\newcounter{dummy} 
\numberwithin{dummy}{section}
\theoremstyle{ocrenumbox}
\newtheorem{theoremeT}[dummy]{Theorem}
\newtheorem{problem}{Problem}[chapter]
\newtheorem{exerciseT}{Exercise}[chapter]
\theoremstyle{blacknumex}
\newtheorem{exampleT}{Example}[chapter]
\theoremstyle{blacknumbox}
\newtheorem{vocabulary}{Vocabulary}[chapter]
\newtheorem{definitionT}{Definition}[section]
\newtheorem{corollaryT}[dummy]{Corollary}
\theoremstyle{ocrenum}
\newtheorem{proposition}[dummy]{Proposition}

\theoremstyle{coloredbox}
\newtheorem{quadroT}{Quadro}[chapter]
%\newtheorem{quadroT}{Box}[chapter] %english

%----------------------------------------------------------------------------------------
%	DEFINITION OF COLORED BOXES
%----------------------------------------------------------------------------------------

\RequirePackage[framemethod=default]{mdframed} % Required for creating the theorem, definition, exercise and corollary boxes

% Theorem box
\newmdenv[skipabove=7pt,
skipbelow=7pt,
backgroundcolor=black!5,
linecolor=color,
innerleftmargin=5pt,
innerrightmargin=5pt,
innertopmargin=5pt,
leftmargin=0cm,
rightmargin=0cm,
innerbottommargin=5pt]{tBox}

% Exercise box	  
\newmdenv[skipabove=7pt,
skipbelow=7pt,
rightline=false,
leftline=true,
topline=false,
bottomline=false,
backgroundcolor=color!10,
linecolor=color,
innerleftmargin=5pt,
innerrightmargin=5pt,
innertopmargin=5pt,
innerbottommargin=5pt,
leftmargin=0cm,
rightmargin=0cm,
linewidth=4pt]{eBox}

% Quadro	  
\newmdenv[skipabove=7pt,
skipbelow=7pt,
rightline=false,
leftline=true,
topline=false,
bottomline=false,
backgroundcolor=color!30,
linecolor=color,
innerleftmargin=5pt,
innerrightmargin=5pt,
innertopmargin=5pt,
innerbottommargin=5pt,
leftmargin=0cm,
rightmargin=0cm,
linewidth=4pt]{iBox}

% Quadro branco  
\newmdenv[skipabove=7pt,
skipbelow=7pt,
rightline=true,
leftline=true,
topline=true,
bottomline=true,
backgroundcolor=white,
linecolor=color,
innerleftmargin=15pt,
innerrightmargin=15pt,
innertopmargin=5pt,
innerbottommargin=5pt,
leftmargin=0cm,
rightmargin=0cm,
linewidth=2pt]{clearBox}

% Definition box
\newmdenv[skipabove=7pt,
skipbelow=7pt,
rightline=false,
leftline=true,
topline=false,
bottomline=false,
linecolor=color,
innerleftmargin=5pt,
innerrightmargin=5pt,
innertopmargin=0pt,
leftmargin=0cm,
rightmargin=0cm,
linewidth=4pt,
innerbottommargin=0pt]{dBox}	

% Corollary box
\newmdenv[skipabove=7pt,
skipbelow=7pt,
rightline=false,
leftline=true,
topline=false,
bottomline=false,
linecolor=gray,
backgroundcolor=black!5,
innerleftmargin=5pt,
innerrightmargin=5pt,
innertopmargin=5pt,
leftmargin=0cm,
rightmargin=0cm,
linewidth=4pt,
innerbottommargin=5pt]{cBox}				  
		  

% Creates an environment for each type of theorem and assigns it a theorem text style from the "Theorem Styles" section above and a colored box from above
\newenvironment{theorem}{\begin{tBox}\begin{theoremeT}}{\end{theoremeT}\end{tBox}}
\newenvironment{exercise}{\begin{eBox}\begin{exerciseT}}{\hfill{\color{color}\tiny\ensuremath{\blacksquare}}\end{exerciseT}\end{eBox}}				  
\newenvironment{definition}{\begin{dBox}\begin{definitionT}}{\end{definitionT}\end{dBox}}	
\newenvironment{example}{\begin{exampleT}}{\hfill{\tiny\ensuremath{\blacksquare}}\end{exampleT}}		
\newenvironment{corollary}{\begin{cBox}\begin{corollaryT}}{\end{corollaryT}\end{cBox}}

\newenvironment{infobox}
  {\begin{iBox}\begin{quadroT}}
    {\hfill{\color{color}\tiny\ensuremath{\blacksquare}}
   \end{quadroT}\end{iBox}}  
   
\newenvironment{clearinfobox}
  {\begin{clearBox}\begin{quadroT}}
    {\hfill{\color{color}\tiny\ensuremath{\blacksquare}}
   \end{quadroT}\end{clearBox}}    


%----------------------------------------------------------------------------------------
%	REMARK ENVIRONMENT
%----------------------------------------------------------------------------------------

\newenvironment{remark}{\par\vskip10pt\small % Vertical white space above the remark and smaller font size
\begin{list}{}{
\leftmargin=35pt % Indentation on the left
\rightmargin=25pt}\item\ignorespaces % Indentation on the right
\makebox[-2.5pt]{\begin{tikzpicture}[overlay]
\node[draw=color!60,line width=1pt,circle,fill=color!25,font=\sffamily\bfseries,inner sep=2pt,outer sep=0pt] at (-15pt,0pt){\textcolor{color}{N}};\end{tikzpicture}} % Orange R in a circle
\advance\baselineskip -1pt}{\end{list}\vskip5pt} % Tighter line spacing and white space after remark

%----------------------------------------------------------------------------------------
%	SECTION NUMBERING IN THE MARGIN
%----------------------------------------------------------------------------------------

\makeatletter
\renewcommand{\@seccntformat}[1]{\llap{\textcolor{color}{\csname the#1\endcsname}\hspace{1em}}}                    
\renewcommand{\section}{\@startsection{section}{1}{\z@}
{-4ex \@plus -1ex \@minus -.4ex}
{1ex \@plus.2ex }
{\normalfont\large\sffamily\bfseries}}
\renewcommand{\subsection}{\@startsection {subsection}{2}{\z@}
{-3ex \@plus -0.1ex \@minus -.4ex}
{0.5ex \@plus.2ex }
{\normalfont\sffamily\bfseries}}
\renewcommand{\subsubsection}{\@startsection {subsubsection}{3}{\z@}
{-2ex \@plus -0.1ex \@minus -.2ex}
{0.2ex \@plus.2ex }
{\normalfont\small\sffamily\bfseries}}                        
\renewcommand\paragraph{\@startsection{paragraph}{4}{\z@}
{-2ex \@plus-.2ex \@minus .2ex}
{0.1ex}
{\normalfont\small\sffamily\bfseries}}

%----------------------------------------------------------------------------------------
%	CHAPTER HEADINGS
%----------------------------------------------------------------------------------------

\newcommand{\thechapterimage}{}
\newcommand{\chapterimage}[1]{\renewcommand{\thechapterimage}{#1}}
\def\thechapter{\arabic{chapter}}
\def\@makechapterhead#1{
\thispagestyle{empty}
{\centering \normalfont\sffamily
\ifnum \c@secnumdepth >\m@ne
\if@mainmatter
\startcontents
\begin{tikzpicture}[remember picture,overlay]
\node at (current page.north west)
{\begin{tikzpicture}[remember picture,overlay]

\node[anchor=north west,inner sep=0pt] at (0,0) {\includegraphics[width=\paperwidth]{\thechapterimage}};

%Commenting the 3 lines below removes the small contents box in the chapter heading
\draw[fill=white,opacity=1] (1cm, 0) rectangle (8cm,-7cm);
\node[anchor=north west] at (1cm,.25cm) {\parbox[t][8cm][t]{6.5cm}{\huge\bfseries\flushleft \printcontents{l}{1}{\setcounter{tocdepth}{2}}}};
%\node[anchor=north west] at (1cm,.25cm) {\parbox[t][8cm][t]{6.5cm}{\huge\bfseries\flushleft \printcontents{l}{1}{\setcounter{tocdepth}{1}}}};
\draw[anchor=west] (5cm,-9cm) node [rounded corners=25pt,fill=white,fill opacity=1,text opacity=1,draw=color,draw opacity=1,line width=2pt,inner sep=15pt]{\huge\sffamily\bfseries\textcolor{black}{\thechapter\ ---\ #1\vphantom{plPQq}\makebox[22cm]{}}};

\end{tikzpicture}};
\end{tikzpicture}}\par\vspace*{230\p@}
\fi
\fi
}
\def\@makeschapterhead#1{
\thispagestyle{empty}
{\centering \normalfont\sffamily
\ifnum \c@secnumdepth >\m@ne
\if@mainmatter
%\startcontents
\begin{tikzpicture}[remember picture,overlay]
\node at (current page.north west)
{\begin{tikzpicture}[remember picture,overlay]
%\node[anchor=north west,inner sep=0pt] at (0,0) {\includegraphics[width=\paperwidth]{\thechapterimage}};
%Commenting the 3 lines below removes the small contents box in the chapter heading
%\draw[fill=white,opacity=1] (1cm,0) rectangle (8cm,-7cm);
%\node[anchor=north west] at (1cm,.25cm) {\parbox[t][8cm][t]{6.5cm}{\huge\bfseries\flushleft \printcontents{l}{1}{\setcounter{tocdepth}{2}}}};
\draw[anchor=west] (-0.2cm,-2.5cm) node [fill=white,fill opacity=1,text opacity=1,draw=color,draw opacity=1,line width=2pt,inner sep=15pt, inner xsep=60pt]{\huge\sffamily\bfseries\textcolor{black}{  \ #1\vphantom{plPQq}\makebox[22cm]{}}};
\end{tikzpicture}};
\end{tikzpicture}}\par\vspace*{50\p@}
\fi
\fi
}
\makeatother % Insert the commands.tex file which contains the majority of the structure behind the template

%----------------------------------------------------------------------------------------
\begin{document}

% prints author names as small caps
\renewcommand{\mkbibnamefirst}[1]{\textsc{#1}}
\renewcommand{\mkbibnamelast}[1]{\textsc{#1}}
\renewcommand{\mkbibnameprefix}[1]{\textsc{#1}}
\renewcommand{\mkbibnameaffix}[1]{\textsc{#1}}

\cleardoublepage % Forces the first chapter to start on an odd page so it's on the right

\graphicspath{ {./Images/}}

%----------------------------------------------------------------------------------------
%	TITLE PAGE
%----------------------------------------------------------------------------------------

\thispagestyle{empty}
%\AddToShipoutPicture*{\put(6,5){\includegraphics[scale=1]{background}}} % Image background

\begin{center}
    \vspace*{1cm}
    
    \textsc{\huge\sffamily Your Name}
    
    \vspace*{3cm}
    
    \textsc{\textbf{\huge\sffamily {Your title \\}}}

    \vspace*{2cm}
    
    \textsc{\textbf{\huge\sffamily {Your title in another relevant language\\}}} 
   
    \vspace*{10.6cm}


    \large{\sffamily{Your city}\\
    \sffamily{\the\year}}
\end{center}


%----------------------------------------------------------------------------------------
%	COPYRIGHT PAGE - TO BE USED IN BOOKS.
%       Most likely not applicable to thesis
%----------------------------------------------------------------------------------------

\iffalse
\newpage
~\vfill
\thispagestyle{empty}

\noindent Copyright \copyright\ \the\year Your Name\\ % Copyright notice

\noindent \textsc{Published by Publisher}\\ % Publisher

\noindent \textsc{book-website.com}\\ % URL

\noindent Licensed under the Creative Commons Attribution-NonCommercial 3.0 Unported License (the ``License''). You may not use this file except in compliance with the License. You may obtain a copy of the License at \url{http://creativecommons.org/licenses/by-nc/3.0}. Unless required by applicable law or agreed to in writing, software distributed under the License is distributed on an \textsc{``AS IS'' BASIS, WITHOUT WARRANTIES OR CONDITIONS OF ANY KIND}, either express or implied. See the License for the specific language governing permissions and limitations under the License.\\ % License information

\noindent \textit{First printing, some date} % Printing/edition date
\fi

%----------------------------------------------------------------------------------------
%	PÁGINA DE ROSTO
%----------------------------------------------------------------------------------------

\newpage
\thispagestyle{empty}
\begin{center}
    \vspace*{1cm}
    
    \textsc{\huge\sffamily Seu Nome}
    
    \vspace*{3cm}
    
    \textsc{\textbf{\huge\sffamily {Seu título \\}}}

    \vspace*{2cm}
    
    \textsc{\textbf{\huge\sffamily {Seu título em inglês \\}}}
    
    \vspace*{10cm} %distance to location and date
    
    %\iffalse
    \begin{tikzpicture}[remember picture,overlay]
     \node[white,draw,text=black,text width=8.5cm] at (page cs:0.4,-0.3) 
     {\sffamily\textbf Versão original da tese apresentada ao <instituto em questão> da <universidade em questão>
     para obtenção do título de Doutor em Ciências, na <aárea em questão>\\
     \vspace{\baselineskip}
     Orientador: Prof. Dr. };
    \end{tikzpicture}
    %\fi

    \iffalse
    \begin{tikzpicture}[remember picture,overlay]
     \node[white,draw,text=black,text width=8.5cm] at (page cs:0.4,-0.3) 
     {\sffamily\textbf Esta versão da tese contém as correções e alterações sugeridas
      pela comissão julgadora durante a defesa da versão original do trabalho,
      realizada em <data da defesa>. Uma cópia da versão original está disponível no
      <instituto em questão> da <universidade em questão>.\\
     \vspace{\baselineskip}
     Orientador: Prof. Dr. };
    \end{tikzpicture}
    \fi

    \large{\sffamily{São Paulo}\\
    \sffamily{\the\year}}
\end{center}

%----------------------------------------------------------------------------------------
%	FICHA CATALOGRÁFICA
%----------------------------------------------------------------------------------------

\newpage
\thispagestyle{empty}

\begin{center}

{\Large\sffamily\textbf {Ficha catalográfica}}

\vspace*{11 cm}

\begin{tikzpicture}[remember picture,overlay]
 \node[draw,text width=9cm] at (page cs:0,0.48) 
 {\sffamily\large {Sobrenome, Nome\\
  {\indent} Título em português\\
  {\indent} Número de páginas: \pageref{LastPage} páginas\\
  \vspace{\baselineskip}
  {\indent} Tese (Doutorado) - <instituto em questão> da < universidade em questão>.
  <departamento em questão>.\\
  \vspace{\baselineskip}
  {\indent} 1. Primeira palavra-chave 2. Segunda palavra-chave 3. Terceira palavra-chave. 
  I. <universidade em questão>. <instituto em questão>.
  <departamento em questão>.\\}
 };
\end{tikzpicture}


{\LARGE\sffamily {Comissão Julgadora:}}

\vspace{1.5cm}
\begin{tabular*}{\textwidth}{@{\extracolsep{\fill}}l l}
\rule{16em}{1px} 	& \rule{16em}{1px} \\
Prof. Dr. 		& Prof. Dr.
\end{tabular*}

\vspace{1.5cm}
\begin{tabular*}{\textwidth}{@{\extracolsep{\fill}}l l}
\rule{16em}{1px} 	& \rule{16em}{1px} \\
Prof. Dr. 		& Prof. Dr.
\end{tabular*}

\vspace{1.5cm}
\begin{tabular*}{\textwidth}{@{\extracolsep{\fill}}l l}
\rule{16em}{1px} 	& \\
Prof. Dr. 	& \\
(Orientador) &
\end{tabular*}

\end{center}
\pagebreak


%----------------------------------------------------------------------------------------
%	Foreword
%----------------------------------------------------------------------------------------
\chapter*{Foreword}
\thispagestyle{plain}

This is a sample of the LaTeX template which was used on my graduate dissertation. It is
heavily based on The Legrand Orange Book template created by Mathias Legrand and 
modified by Vel. The original template was downloaded from 
\url{http://www.latextemplates.com/template/the-legrand-orange-book} and subsequently modified
to suit my personal taste and needs. The format and style do not follow the guidelines of any
particular university.

In this sample, I will outline the main features of this template and
point out some missing features (Chapter \ref{chap:main}), exemplify how it can be used when
the dissertation chapter corresponds to a research article (Chapter \ref{chap:article}) and 
outline additional features
of the Legrand template which are still available in the current template, but may be less
applicable to dissertations (Chapter \ref{chap:legrand}). This sample, however, should
NOT be taken as the template documentation! Most information required to use the template
is writen on the ``thesis.tex'' and ``structure.tex'' files.

\vspace{1.5cm}
\begin{flushright}
 \textit{Ayana B. Martins}
\end{flushright}




%----------------------------------------------------------------------------------------
%	AGRADECIMENTOS
%----------------------------------------------------------------------------------------
\chapter*{Agradecimentos}
\thispagestyle{plain}

\lipsum[1-3]


%----------------------------------------------------------------------------------------
%	RESUMO
%----------------------------------------------------------------------------------------
\chapter*{Resumo}
\thispagestyle{plain}

\lipsum[1-3]



%----------------------------------------------------------------------------------------
%	ABSTRACT
%----------------------------------------------------------------------------------------
\chapter*{Abstract}
\thispagestyle{plain}

\lipsum[1-3]


%----------------------------------------------------------------------------------------
%	TABLE OF CONTENTS
%----------------------------------------------------------------------------------------
\tableofcontents % Print the table of contents itself
\thispagestyle{plain} % No headers

\cleardoublepage % Forces the first chapter to start on an odd page so it's on the right

\pagestyle{fancy} % Print headers again



\mainmatter

%Each chapter can be entered as a different file, which can be compiled separately with
%few changes

%----------------------------------------------------------------------------------------
%	CAPÍTULO 1
%----------------------------------------------------------------------------------------
%To compile this chapter separately, comment out the \iffalse and \fi lines

\iffalse
\documentclass[11pt,fleqn]{book} % Default font size and left-justified equations
\usepackage[top=3cm,bottom=3cm,left=3.2cm,right=3.2cm,headsep=10pt,a4paper]{geometry} % Page margins
\usepackage{xcolor} % Required for specifying colors by name
\definecolor{color}{RGB}{128,0,0} % Define the color used for highlighting throughout the book
\definecolor{blue}{RGB}{69,74,159} %Redifines the default for the color blue

%----------------------------------------------------------------------------------------
% Font Settings
\usepackage{avant} % Use the Avantgarde font for headings
%\usepackage{times} % Use the Times font for headings
\usepackage{mathptmx} % Use the Adobe Times Roman as the default text font together with math symbols from the Sym­bol, Chancery and Com­puter Modern fonts
\usepackage{microtype} % Slightly tweak font spacing for aesthetics
\usepackage[utf8]{inputenc} % Required for including letters with accents
\usepackage[T1]{fontenc} % Use 8-bit encoding that has 256 glyphs
\usepackage{csquotes}
\usepackage{endnotes}
\usepackage{setspace}
\onehalfspacing

\usepackage{todonotes}

%----------------------------------------------------------------------------------------
% Bibliography
\usepackage[style=authoryear,
            %sorting=nyt,
            %sortcites=true,
            autopunct=true,
            %babel=hyphen,
            %hyperref=true,
            %abbreviate=false,
            backref=true,
            %refsection=chapter,
            backend=bibtex]
{biblatex}
\addbibresource{../../tese.bib} % BibTeX bibliography file
\defbibheading{bibempty}{}
%bibtex sspanel1-blx

%Another option with natbib which require a different method to do a bibliography section per chapter
%\usepackage[square,numbers,sectionbib]{natbib}
%\usepackage{chapterbib}
%\usepackage{authblk}
%\bibliographystyle{evolution}

%----------------------------------------------------------------------------------------
% Index
\usepackage{calc} % For simpler calculation - used for spacing the index letter headings correctly
\usepackage{makeidx} % Required to make an index
\makeindex % Tells LaTeX to create the files required for indexing

%----------------------------------------------------------------------------------------
%----------------------------------------------------------------------------------------
%	VARIOUS REQUIRED PACKAGES
%----------------------------------------------------------------------------------------

\usepackage{titlesec} % Allows customization of titles
\usepackage{graphicx} % Required for including pictures
\graphicspath{{Images/}} % Specifies the directory where pictures are stored
\usepackage{lipsum} % Inserts dummy text
\usepackage{tikz} % Required for drawing custom shapes
\usepackage[brazil]{babel} % Specify language/hyphenation
\usepackage{enumitem} % Customize lists
\setlist{nolistsep} % Reduce spacing between bullet points and numbered lists
\usepackage{booktabs} % Required for nicer horizontal rules in tables
\usepackage{eso-pic} % Required for specifying an image background in the title page
\usepackage[rightcaption]{sidecap} % Required for placing caption on the side of float 

%----------------------------------------------------------------------------------------
%	MAIN TABLE OF CONTENTS
%----------------------------------------------------------------------------------------

\usepackage{titletoc} % Required for manipulating the table of contents

\addtocontents{toc}{\protect\thispagestyle{empty}}
\contentsmargin{0cm} % Removes the default margin
% Chapter text styling
\titlecontents{chapter}[1.25cm] % Indentation
{\addvspace{15pt}\large\sffamily\bfseries} % Spacing and font options for chapters
{\color{color!60}\contentslabel[\Large\thecontentslabel]{1.25cm}\color{color}} % Chapter number
{}  
{\color{color!60}\normalsize\sffamily\bfseries\;\titlerule*[.5pc]{.}\;\thecontentspage} % Page number
% Section text styling
\titlecontents{section}[1.25cm] % Indentation
{\addvspace{5pt}\sffamily\bfseries} % Spacing and font options for sections
{\contentslabel[\thecontentslabel]{1.25cm}} % Section number
{}
{\sffamily\hfill\color{black}\thecontentspage} % Page number
[]
% Subsection text styling
\titlecontents{subsection}[1.25cm] % Indentation
{\addvspace{1pt}\sffamily\small} % Spacing and font options for subsections
{\contentslabel[\thecontentslabel]{1.25cm}} % Subsection number
{}
{\sffamily\;\titlerule*[.5pc]{.}\;\thecontentspage} % Page number
[] 

%----------------------------------------------------------------------------------------
%	MINI TABLE OF CONTENTS IN CHAPTER HEADS
%----------------------------------------------------------------------------------------

% Section text styling
\titlecontents{lsection}[0em] % Indendating
{\footnotesize\sffamily} % Font settings
{}
{}
{}

% Subsection text styling
\titlecontents{lsubsection}[.5em] % Indentation
{\normalfont\footnotesize\sffamily} % Font settings
{}
{}
{}
 
%----------------------------------------------------------------------------------------
%	PAGE HEADERS
%----------------------------------------------------------------------------------------

\usepackage{fancyhdr} % Required for header and footer configuration

\pagestyle{fancy}
\pagenumbering{roman} 
\renewcommand{\chaptermark}[1]{\markboth{\sffamily\normalsize\bfseries #1}{}} % Chapter text font settings
\renewcommand{\sectionmark}[1]{\markright{\sffamily\normalsize\thesection\hspace{5pt}#1}{}} % Section text font settings


%-----------------------------------------------------------------------------------------
%Two-sided printing
%\fancyhf{} \fancyhead[LE,RO]{\sffamily\normalsize\thepage} % Font setting for the page number in the header
%\fancyhead[LO]{\rightmark} % Print the nearest section name on the left side of odd pages
%\fancyhead[RE]{\leftmark} % Print the current chapter name on the right side of even pages
%------------------------------------------------------------------------------------------

%-----------------------------------------------------------------------------------------
%One-sided printing
\fancyhf{} \fancyhead[RE,RO]{\sffamily\normalsize\thepage} % Font setting and placement for the page number in the header
\fancyhead[LO]{\rightmark} % Print the nearest section name on the left side of odd pages
\fancyhead[LE]{\leftmark} % Print the current chapter name on the right side of even pages
%------------------------------------------------------------------------------------------

\renewcommand{\headrulewidth}{0.5pt} % Width of the rule under the header
\addtolength{\headheight}{2.5pt} % Increase the spacing around the header slightly
\renewcommand{\footrulewidth}{0pt} % Removes the rule in the footer
\fancypagestyle{plain}{\fancyfoot[CO,CE]{\thepage}
                      \renewcommand{\headrulewidth}{0pt}} % Style for when a plain pagestyle is specified


\makeatletter
% Removes the header from odd empty pages at the end of chapters
%\renewcommand{\cleardoublepage}{
%\clearpage\ifodd\c@page\else
%\hbox{}
%\vspace*{\fill}
%\thispagestyle{empty}
%\newpage
%\fi}

%----------------------------------------------------------------------------------------
%	SET UP POSITION OF TEXT IN TITLE PAGE
%----------------------------------------------------------------------------------------
\def\parsecomma#1,#2\endparsecomma{\def\page@x{#1}\def\page@y{#2}}
\tikzdeclarecoordinatesystem{page}{
    \parsecomma#1\endparsecomma
    \pgfpointanchor{current page}{north east}
    % Save the upper right corner
    \pgf@xc=\pgf@x%
    \pgf@yc=\pgf@y%
    % save the lower left corner
    \pgfpointanchor{current page}{south west}
    \pgf@xb=\pgf@x%
    \pgf@yb=\pgf@y%
    % Transform to the correct placement
    \pgfmathparse{(\pgf@xc-\pgf@xb)/2.*\page@x+(\pgf@xc+\pgf@xb)/2.}
    \expandafter\pgf@x\expandafter=\pgfmathresult pt
    \pgfmathparse{(\pgf@yc-\pgf@yb)/2.*\page@y+(\pgf@yc+\pgf@yb)/2.}
    \expandafter\pgf@y\expandafter=\pgfmathresult pt
}

%----------------------------------------------------------------------------------------
%	MODIFIED ENDNOTE AND FOOTNOTES COMMANDS
%----------------------------------------------------------------------------------------
\renewcommand{\theendnote}{\alph{endnote}}

\newcommand\footnoteref[1]{\protected@xdef\@thefnmark{\ref{#1}}\@footnotemark}

%Use the following command to add a endnote to a footnote without showing the number.

\newcommand\numberlessnote[1]{%
  \begingroup
  \renewcommand\thefootnote{}\footnote{#1}%
  \addtocounter{footnote}{-1}%
  \endgroup
}

%\newcommand\insidenote[1]{%
%  \begingroup
%  \addtocounter{endnote}{+1}%
%  \renewcommand\theendnotes{}\endnotetext{#1}%
%  \endgroup
%}

%Use the following command if you are using insidenotes in our text and one to keep which
%note number unique
%Example:
% Some text \footnote{See page \ref{endnotes} \insidenote{Actual note}}
% The label endnote should then be placed in the page with the notes.

%\newcommand\onlyfootnote[1]{%
%  \begingroup
%  \addtocounter{endnote}{+1}%
%  \renewcommand\theendnotes{}\footnote{#1}%
%  \endgroup
%}

%\renewcommand*{\thefootnote}{\fnsymbol{footnote}}

%----------------------------------------------------------------------------------------
%	MODIFIED TABLES
%----------------------------------------------------------------------------------------

\usepackage{tabularx}
\newcolumntype{C}[1]{>{\hsize=#1\hsize\centering\arraybackslash}X}%
\newcolumntype{L}[1]{>{\hsize=#1\hsize\raggedright\arraybackslash}X}%


\def\hlinewd#1{%
\noalign{\ifnum0=`}\fi\hrule \@height #1 %
\futurelet\reserved@a\@xhline}


\newcommand{\abstractseparator}{
\textcolor{color}{\rule[0.5ex]{1\columnwidth}{2pt}}
}

%----------------------------------------------------------------------------------------
%	THEOREM STYLES
%----------------------------------------------------------------------------------------


\usepackage{amsmath,amsfonts,amssymb,amsthm} % For including math equations, theorems, symbols, etc

\newcommand{\intoo}[2]{\mathopen{]}#1\,;#2\mathclose{[}}
\newcommand{\ud}{\mathop{\mathrm{{}d}}\mathopen{}}
\newcommand{\intff}[2]{\mathopen{[}#1\,;#2\mathclose{]}}
\newtheorem{notation}{Notation}[chapter]

%%%%%%%%%%%%%%%%%%%%%%%%%%%%%%%%%%%%%%%%%%%%%%%%%%%%%%%%%%%%%%%%%%%%%%%%%%%
%%%%%%%%%%%%%%%%%%%% dedicated to boxed/framed environements %%%%%%%%%%%%%%
%%%%%%%%%%%%%%%%%%%%%%%%%%%%%%%%%%%%%%%%%%%%%%%%%%%%%%%%%%%%%%%%%%%%%%%%%%%

\newtheoremstyle{ocrenumbox}% % Theorem style name
{0pt}% Space above
{0pt}% Space below
{\normalfont}% % Body font
{}% Indent amount
{\small\bf\sffamily\color{color}}% % Theorem head font
{\;}% Punctuation after theorem head
{0.25em}% Space after theorem head
{\small\sffamily\color{color}\thmname{#1}\nobreakspace\thmnumber{\@ifnotempty{#1}{}\@upn{#2}}% Theorem text (e.g. Theorem 2.1)
\thmnote{\nobreakspace\the\thm@notefont\sffamily\bfseries\color{black}---\nobreakspace#3.}} % Optional theorem note
\renewcommand{\qedsymbol}{$\blacksquare$}% Optional qed square

\newtheoremstyle{blacknumex}% Theorem style name
{5pt}% Space above
{5pt}% Space below
{\normalfont}% Body font
{} % Indent amount
{\small\bf\sffamily}% Theorem head font
{\;}% Punctuation after theorem head
{0.25em}% Space after theorem head
{\small\sffamily{\tiny\ensuremath{\blacksquare}}\nobreakspace\thmname{#1}\nobreakspace\thmnumber{\@ifnotempty{#1}{}\@upn{#2}}% Theorem text (e.g. Theorem 2.1)
\thmnote{\nobreakspace\the\thm@notefont\sffamily\bfseries---\nobreakspace#3.}}% Optional theorem note

\newtheoremstyle{blacknumbox} % Theorem style name
{0pt}% Space above
{0pt}% Space below
{\normalfont}% Body font
{}% Indent amount
{\small\bf\sffamily}% Theorem head font
{\;}% Punctuation after theorem head
{0.25em}% Space after theorem head
{\small\sffamily\thmname{#1}\nobreakspace\thmnumber{\@ifnotempty{#1}{}\@upn{#2}}% Theorem text (e.g. Theorem 2.1)
\thmnote{\nobreakspace\the\thm@notefont\sffamily\bfseries---\nobreakspace#3.}}% Optional theorem note


\newtheoremstyle{coloredbox}% Box style name
{0pt}% Space above
{0pt}% Space below
{\normalfont}% % Body font
{}% Indent amount
{\small\bf\sffamily\color{color}}% % Theorem head font
{\;}% Punctuation after theorem head
{\newline}% Space after theorem head
{\small\sffamily\color{color}\thmname{#1}\nobreakspace\thmnumber{\@ifnotempty{#1}{}\@upn{#2}}% Theorem text (e.g. Theorem 2.1)
\thmnote{\nobreakspace\the\thm@notefont\sffamily\bfseries\color{black}---\nobreakspace#3.}} % Optional theorem note
\renewcommand{\qedsymbol}{$\blacksquare$}% Optional qed square

%%%%%%%%%%%%%%%%%%%%%%%%%%%%%%%%%%%%%%%%%%%%%%%%%%%%%%%%%%%%%%%%%%%%%%%%%%%
%%%%%%%%%%%%% dedicated to non-boxed/non-framed environements %%%%%%%%%%%%%
%%%%%%%%%%%%%%%%%%%%%%%%%%%%%%%%%%%%%%%%%%%%%%%%%%%%%%%%%%%%%%%%%%%%%%%%%%%
\newtheoremstyle{ocrenum}% % Theorem style name
{5pt}% Space above
{5pt}% Space below
{\normalfont}% % Body font
{}% Indent amount
{\small\bf\sffamily\color{color}}% % Theorem head font
{\;}% Punctuation after theorem head
{0.25em}% Space after theorem head
{\small\sffamily\color{color}\thmname{#1}\nobreakspace\thmnumber{\@ifnotempty{#1}{}\@upn{#2}}% Theorem text (e.g. Theorem 2.1)
\thmnote{\nobreakspace\the\thm@notefont\sffamily\bfseries\color{black}---\nobreakspace#3.}} % Optional theorem note
\renewcommand{\qedsymbol}{$\blacksquare$}% Optional qed square
\makeatother

% Defines the theorem text style for each type of theorem to one of the three styles above
\newcounter{dummy} 
\numberwithin{dummy}{section}
\theoremstyle{ocrenumbox}
\newtheorem{theoremeT}[dummy]{Theorem}
\newtheorem{problem}{Problem}[chapter]
\newtheorem{exerciseT}{Exercise}[chapter]
\theoremstyle{blacknumex}
\newtheorem{exampleT}{Example}[chapter]
\theoremstyle{blacknumbox}
\newtheorem{vocabulary}{Vocabulary}[chapter]
\newtheorem{definitionT}{Definition}[section]
\newtheorem{corollaryT}[dummy]{Corollary}
\theoremstyle{ocrenum}
\newtheorem{proposition}[dummy]{Proposition}

\theoremstyle{coloredbox}
\newtheorem{quadroT}{Quadro}[chapter]
%\newtheorem{quadroT}{Box}[chapter] %english

%----------------------------------------------------------------------------------------
%	DEFINITION OF COLORED BOXES
%----------------------------------------------------------------------------------------

\RequirePackage[framemethod=default]{mdframed} % Required for creating the theorem, definition, exercise and corollary boxes

% Theorem box
\newmdenv[skipabove=7pt,
skipbelow=7pt,
backgroundcolor=black!5,
linecolor=color,
innerleftmargin=5pt,
innerrightmargin=5pt,
innertopmargin=5pt,
leftmargin=0cm,
rightmargin=0cm,
innerbottommargin=5pt]{tBox}

% Exercise box	  
\newmdenv[skipabove=7pt,
skipbelow=7pt,
rightline=false,
leftline=true,
topline=false,
bottomline=false,
backgroundcolor=color!10,
linecolor=color,
innerleftmargin=5pt,
innerrightmargin=5pt,
innertopmargin=5pt,
innerbottommargin=5pt,
leftmargin=0cm,
rightmargin=0cm,
linewidth=4pt]{eBox}

% Quadro	  
\newmdenv[skipabove=7pt,
skipbelow=7pt,
rightline=false,
leftline=true,
topline=false,
bottomline=false,
backgroundcolor=color!30,
linecolor=color,
innerleftmargin=5pt,
innerrightmargin=5pt,
innertopmargin=5pt,
innerbottommargin=5pt,
leftmargin=0cm,
rightmargin=0cm,
linewidth=4pt]{iBox}

% Quadro branco  
\newmdenv[skipabove=7pt,
skipbelow=7pt,
rightline=true,
leftline=true,
topline=true,
bottomline=true,
backgroundcolor=white,
linecolor=color,
innerleftmargin=15pt,
innerrightmargin=15pt,
innertopmargin=5pt,
innerbottommargin=5pt,
leftmargin=0cm,
rightmargin=0cm,
linewidth=2pt]{clearBox}

% Definition box
\newmdenv[skipabove=7pt,
skipbelow=7pt,
rightline=false,
leftline=true,
topline=false,
bottomline=false,
linecolor=color,
innerleftmargin=5pt,
innerrightmargin=5pt,
innertopmargin=0pt,
leftmargin=0cm,
rightmargin=0cm,
linewidth=4pt,
innerbottommargin=0pt]{dBox}	

% Corollary box
\newmdenv[skipabove=7pt,
skipbelow=7pt,
rightline=false,
leftline=true,
topline=false,
bottomline=false,
linecolor=gray,
backgroundcolor=black!5,
innerleftmargin=5pt,
innerrightmargin=5pt,
innertopmargin=5pt,
leftmargin=0cm,
rightmargin=0cm,
linewidth=4pt,
innerbottommargin=5pt]{cBox}				  
		  

% Creates an environment for each type of theorem and assigns it a theorem text style from the "Theorem Styles" section above and a colored box from above
\newenvironment{theorem}{\begin{tBox}\begin{theoremeT}}{\end{theoremeT}\end{tBox}}
\newenvironment{exercise}{\begin{eBox}\begin{exerciseT}}{\hfill{\color{color}\tiny\ensuremath{\blacksquare}}\end{exerciseT}\end{eBox}}				  
\newenvironment{definition}{\begin{dBox}\begin{definitionT}}{\end{definitionT}\end{dBox}}	
\newenvironment{example}{\begin{exampleT}}{\hfill{\tiny\ensuremath{\blacksquare}}\end{exampleT}}		
\newenvironment{corollary}{\begin{cBox}\begin{corollaryT}}{\end{corollaryT}\end{cBox}}

\newenvironment{infobox}
  {\begin{iBox}\begin{quadroT}}
    {\hfill{\color{color}\tiny\ensuremath{\blacksquare}}
   \end{quadroT}\end{iBox}}  
   
\newenvironment{clearinfobox}
  {\begin{clearBox}\begin{quadroT}}
    {\hfill{\color{color}\tiny\ensuremath{\blacksquare}}
   \end{quadroT}\end{clearBox}}    


%----------------------------------------------------------------------------------------
%	REMARK ENVIRONMENT
%----------------------------------------------------------------------------------------

\newenvironment{remark}{\par\vskip10pt\small % Vertical white space above the remark and smaller font size
\begin{list}{}{
\leftmargin=35pt % Indentation on the left
\rightmargin=25pt}\item\ignorespaces % Indentation on the right
\makebox[-2.5pt]{\begin{tikzpicture}[overlay]
\node[draw=color!60,line width=1pt,circle,fill=color!25,font=\sffamily\bfseries,inner sep=2pt,outer sep=0pt] at (-15pt,0pt){\textcolor{color}{N}};\end{tikzpicture}} % Orange R in a circle
\advance\baselineskip -1pt}{\end{list}\vskip5pt} % Tighter line spacing and white space after remark

%----------------------------------------------------------------------------------------
%	SECTION NUMBERING IN THE MARGIN
%----------------------------------------------------------------------------------------

\makeatletter
\renewcommand{\@seccntformat}[1]{\llap{\textcolor{color}{\csname the#1\endcsname}\hspace{1em}}}                    
\renewcommand{\section}{\@startsection{section}{1}{\z@}
{-4ex \@plus -1ex \@minus -.4ex}
{1ex \@plus.2ex }
{\normalfont\large\sffamily\bfseries}}
\renewcommand{\subsection}{\@startsection {subsection}{2}{\z@}
{-3ex \@plus -0.1ex \@minus -.4ex}
{0.5ex \@plus.2ex }
{\normalfont\sffamily\bfseries}}
\renewcommand{\subsubsection}{\@startsection {subsubsection}{3}{\z@}
{-2ex \@plus -0.1ex \@minus -.2ex}
{0.2ex \@plus.2ex }
{\normalfont\small\sffamily\bfseries}}                        
\renewcommand\paragraph{\@startsection{paragraph}{4}{\z@}
{-2ex \@plus-.2ex \@minus .2ex}
{0.1ex}
{\normalfont\small\sffamily\bfseries}}

%----------------------------------------------------------------------------------------
%	CHAPTER HEADINGS
%----------------------------------------------------------------------------------------

\newcommand{\thechapterimage}{}
\newcommand{\chapterimage}[1]{\renewcommand{\thechapterimage}{#1}}
\def\thechapter{\arabic{chapter}}
\def\@makechapterhead#1{
\thispagestyle{empty}
{\centering \normalfont\sffamily
\ifnum \c@secnumdepth >\m@ne
\if@mainmatter
\startcontents
\begin{tikzpicture}[remember picture,overlay]
\node at (current page.north west)
{\begin{tikzpicture}[remember picture,overlay]

\node[anchor=north west,inner sep=0pt] at (0,0) {\includegraphics[width=\paperwidth]{\thechapterimage}};

%Commenting the 3 lines below removes the small contents box in the chapter heading
\draw[fill=white,opacity=1] (1cm, 0) rectangle (8cm,-7cm);
\node[anchor=north west] at (1cm,.25cm) {\parbox[t][8cm][t]{6.5cm}{\huge\bfseries\flushleft \printcontents{l}{1}{\setcounter{tocdepth}{2}}}};
%\node[anchor=north west] at (1cm,.25cm) {\parbox[t][8cm][t]{6.5cm}{\huge\bfseries\flushleft \printcontents{l}{1}{\setcounter{tocdepth}{1}}}};
\draw[anchor=west] (5cm,-9cm) node [rounded corners=25pt,fill=white,fill opacity=1,text opacity=1,draw=color,draw opacity=1,line width=2pt,inner sep=15pt]{\huge\sffamily\bfseries\textcolor{black}{\thechapter\ ---\ #1\vphantom{plPQq}\makebox[22cm]{}}};

\end{tikzpicture}};
\end{tikzpicture}}\par\vspace*{230\p@}
\fi
\fi
}
\def\@makeschapterhead#1{
\thispagestyle{empty}
{\centering \normalfont\sffamily
\ifnum \c@secnumdepth >\m@ne
\if@mainmatter
%\startcontents
\begin{tikzpicture}[remember picture,overlay]
\node at (current page.north west)
{\begin{tikzpicture}[remember picture,overlay]
%\node[anchor=north west,inner sep=0pt] at (0,0) {\includegraphics[width=\paperwidth]{\thechapterimage}};
%Commenting the 3 lines below removes the small contents box in the chapter heading
%\draw[fill=white,opacity=1] (1cm,0) rectangle (8cm,-7cm);
%\node[anchor=north west] at (1cm,.25cm) {\parbox[t][8cm][t]{6.5cm}{\huge\bfseries\flushleft \printcontents{l}{1}{\setcounter{tocdepth}{2}}}};
\draw[anchor=west] (-0.2cm,-2.5cm) node [fill=white,fill opacity=1,text opacity=1,draw=color,draw opacity=1,line width=2pt,inner sep=15pt, inner xsep=60pt]{\huge\sffamily\bfseries\textcolor{black}{  \ #1\vphantom{plPQq}\makebox[22cm]{}}};
\end{tikzpicture}};
\end{tikzpicture}}\par\vspace*{50\p@}
\fi
\fi
}
\makeatother % Insert the commands.tex file which contains the majority of the structure behind the template

%----------------------------------------------------------------------------------------

\begin{document}
\graphicspath{ {../../Images/}}
\fi


\chapterimage{Headers/chapter_header.pdf}
\chapter{Regular chapter \label{chap:main}}
%\begin{refsection}
%------------------------------------------------
%\listoftodos

\section{Mini table of content in chapter header}
This template includes a small TOC in each chapter. Note that it is NOT formatted so that the content
automatically fits in the box (see chapter \ref{chap:legrand}).


\section{Sections}
\subsection{Subsections}

\begin{itemize}

\item \textbf {\sffamily \footnotesize List of items with text:}
 
Here you can write something about this topic. 

\hspace{15pt} Indentation must be forced if there is a second paragraph within an item.

\item \textbf {\sffamily \footnotesize Another item:}

\lipsum[1]

\end{itemize}

\section{Using citations}

This template uses the package biblatex (\url{https://www.ctan.org/pkg/biblatex})
for reference management. Please mind the compilation instructions in
the main file (``thesis.tex'').

A statement may be followed by single (I love natural selection \parencite{darwin_origin_1859}) or
multiple citations (We love natural selection \parencite{darwin_tendency_1858, darwin_origin_1859}).

References may also appear within a sentence: ``We thank \textcite{legrand_legrand_2013} for much of
this template''. 

Biblatex allows easy customization of parenthetical citations:

\begin{itemize}
  \item \parencite[p. 180]{darwin_origin_1859}.
   \item \parencite[revisão em][]{darwin_origin_1859}.
  \item \parencite[ver][p. 180]{darwin_origin_1859}.
\end{itemize}

\section{Using figures}


The structure.tex file calls the package sidecap (\url{http://www.ctan.org/tex-archive/macros/latex/contrib/sidecap})
which allows captions to be typeset sideways.
The graphic path is defined as ``Images/'' in this same file.


 \begin{figure}
\centering
\includegraphics[width=14cm]{Figures/1-regular-chapter/placeholder.jpg}
 \caption{Regular figure caption}
\end{figure}

\begin{SCfigure}
\centering
\includegraphics[width=6cm]{Figures/1-regular-chapter/placeholder.jpg}
 \caption{
Caption may appear on the side of floats.
 \vspace{8pt}}
\end{SCfigure}


\section{Using tables}

The structure.tex file calls the package tabularx (\url{https://www.ctan.org/pkg/tabularx}) that allows more flexibility in
choosing column widths.


\begin{table}
\centering
\caption{Table caption
}
\begin{tabularx}{\textwidth}{ C{1} | C{1} C{1} C{1} C{1} C{1} C{1} C{1} C{1} C{1} C{1}}
   \hlinewd{1.5pt}
      & \textbf{\textcolor{blue}{A}} & \textbf{\textcolor{blue}{B}} & \textbf{\textcolor{blue}{C}} & \textbf{\textcolor{blue}{D}}
      & \textbf{\textcolor{blue}{E}} & \textbf{\textcolor{blue}{F}} & \textbf{\textcolor{blue}{G}} & \textbf{\textcolor{blue}{H}}
      & \textbf{\textcolor{blue}{I}} & \textbf{\textcolor{blue}{J}}\tabularnewline
   \hline
    \textbf{A} & \textcolor{blue}{7.73/} \newline 5.76 & \vspace{0pt}\textcolor{blue}{7.04} &\vspace{0pt}\textcolor{blue} {7.60}
    & \vspace{0pt} \textcolor{blue}{10.33} &\vspace{0pt}\textcolor{blue}{11.05} & \vspace{0pt} \textcolor{blue}{14.50} 
    &\vspace{0pt} \textcolor{blue}{13.17} & \vspace{0pt}\textcolor{blue}{17.00} &\vspace{0pt} \textcolor{blue}{17.36}
    &\vspace{0pt} \textcolor{blue}{15.78}\tabularnewline
    
    \textbf{B} & \vspace{0pt}6.58 &\textcolor{blue}{6.68/} \newline 6.60 &\vspace{0pt}\textcolor{blue} {7.20}
    & \vspace{0pt} \textcolor{blue}{10.16} &\vspace{0pt}\textcolor{blue}{11.41} & \vspace{0pt} \textcolor{blue}{14.94} 
    &\vspace{0pt} \textcolor{blue}{13.58} & \vspace{0pt}\textcolor{blue}{17.43} &\vspace{0pt} \textcolor{blue}{18.16}
    &\vspace{0pt} \textcolor{blue}{16.04}\tabularnewline
    
    \textbf{C} & \vspace{0pt}8.86 & \vspace{0pt}8.92 &\textcolor{blue}{7.16/} \newline 6.40
    & \vspace{0pt} \textcolor{blue}{10.03} &\vspace{0pt}\textcolor{blue}{10.46} & \vspace{0pt} \textcolor{blue}{13.80} 
    &\vspace{0pt} \textcolor{blue}{12.70} & \vspace{0pt}\textcolor{blue}{15.64} &\vspace{0pt} \textcolor{blue}{16.99}
    &\vspace{0pt} \textcolor{blue}{14.57}\tabularnewline
    
    \textbf{D} & \vspace{0pt}9.52 & \vspace{0pt}9.44 & \vspace{0pt}6.94
    & \textcolor{blue}{10.43/} \newline 6.07 &\vspace{0pt}\textcolor{blue}{10.93} & \vspace{0pt} \textcolor{blue}{14.81} 
    &\vspace{0pt} \textcolor{blue}{13.00} & \vspace{0pt}\textcolor{blue}{14.80} &\vspace{0pt} \textcolor{blue}{16.00}
    &\vspace{0pt} \textcolor{blue}{15.04}\tabularnewline
    
    \textbf{E} & \vspace{0pt}10.16 & \vspace{0pt}10.52 & \vspace{0pt}8.50
    & \vspace{0pt}7.18 &  \textcolor{blue}{10.76/} \newline 6.02 & \vspace{0pt} \textcolor{blue}{14.57} 
    &\vspace{0pt} \textcolor{blue}{12.10} & \vspace{0pt}\textcolor{blue}{14.34} &\vspace{0pt} \textcolor{blue}{16.95}
    &\vspace{0pt} \textcolor{blue}{14.93}\tabularnewline
    
    \textbf{F} & \vspace{0pt}12.56 & \vspace{0pt}13.44 & \vspace{0pt}11.42
    & \vspace{0pt}9.78 &  \vspace{0pt}7.90 & \textcolor{blue}{13.33/} \newline 5.80 
    &\vspace{0pt} \textcolor{blue}{12.42} & \vspace{0pt}\textcolor{blue}{14.10} &\vspace{0pt} \textcolor{blue}{17.00}
    &\vspace{0pt} \textcolor{blue}{15.08}\tabularnewline
    
    \textbf{G} & \vspace{0pt}13.34 & \vspace{0pt}14.62 & \vspace{0pt}12.24
    & \vspace{0pt}10.88 &  \vspace{0pt}8.98 & \vspace{0pt}5.14 
    &\textcolor{blue}{12.27/} \newline 3.93  & \vspace{0pt}\textcolor{blue}{13.00} &\vspace{0pt} \textcolor{blue}{16.44}
    &\vspace{0pt} \textcolor{blue}{14.67}\tabularnewline
    
    \textbf{H} & \vspace{0pt}16.74 & \vspace{0pt}17.54 & \vspace{0pt}14.64
    & \vspace{0pt}13.50 &  \vspace{0pt}10.86 & \vspace{0pt}7.60 
    & \vspace{0pt}6.64   & \textcolor{blue}{10.60/} \newline 5.67 &\vspace{0pt} \textcolor{blue}{14.21}
    &\vspace{0pt} \textcolor{blue}{13.17}\tabularnewline
    
    \textbf{I} & \vspace{0pt}16.74 & \vspace{0pt}17.30 & \vspace{0pt}14.98
    & \vspace{0pt}13.82 &  \vspace{0pt}10.98 & \vspace{0pt}8.10 
    & \vspace{0pt}7.18   & \vspace{0pt}5.98  &\textcolor{blue}{13.64/} \newline 5.31
    &\vspace{0pt} \textcolor{blue}{14.52}\tabularnewline
    
    \textbf{J} & \vspace{0pt}17.36 & \vspace{0pt}8.20 & \vspace{0pt}16.54
    & \vspace{0pt}16.38 &  \vspace{0pt}14.10 & \vspace{0pt}10.46 
    & \vspace{0pt}8.86   & \vspace{0pt}8.18  & \vspace{0pt}7.46
    &\textcolor{blue}{13.40/} \newline 6.47\tabularnewline
    \hlinewd{1.5pt}
\end{tabularx}
\end{table}

\pagebreak

\section{Using boxes}


The ``clearinfobox'' environment can be used to highlight relevant information.
Additionally, this template includes all in-text elements from the Legrand Orange Book
(see chapter \ref{chap:legrand}). The word ``Quadro'' appearing in the box does not change 
automatically when the document language is changed.

\begin{clearinfobox}[Box title]
%Uncomment the following line to make the box appear in the table of contents
%\addcontentsline{toc}{subsection}{\small\sffamily\color{green} Quadro 1.1 - Espécies em anel: conceitos importantes}
%Boxes may include images
 %\includegraphics[width=14cm]{Figures/1-regular-chapter/placeholder.jpg}
 %and text
\lipsum[1]
\end{clearinfobox}

\section{Using footnotes}


There are two types of footnotes:

\begin{itemize}
  \item Short footnotes which appear in the same page \footnote{This is a short footnote}.
  \item Longer footnotes which will appear in a separate page at the end of the document
  and are identified by a letter\footnote{\label{same_number} 
Ver página \pageref{notas}}\endnote {Explanation that will appear in a separate page.}.
\end{itemize}

When there is more than one long footnote in the same page (yes, I do love footnotes) the first
one should be labeled and the additional footnotes should carry the samel label so that they only differ 
in the letter \footnoteref{same_number}\endnote{More info!!}.


%------------------------------------------------
\printbibliography[heading=subbibliography]
%\end{refsection}

\iffalse
\begingroup
\let\enotesize\normalsize
\renewcommand\notesname{Notas finais}
\renewcommand\enoteformat{\noindent[\theenmark]\hspace{0.06in}}
\label{notas}
\theendnotes
\endgroup

\end{document}
\fi



%----------------------------------------------------------------------------------------
%	CAPÍTULO 2
%----------------------------------------------------------------------------------------
\iffalse
\documentclass[11pt,fleqn]{book} % Default font size and left-justified equations
\usepackage[top=3cm,bottom=3cm,left=3.2cm,right=3.2cm,headsep=10pt,a4paper]{geometry} % Page margins
\usepackage{xcolor} % Required for specifying colors by name
\definecolor{color}{RGB}{128,0,0} % Define the color used for highlighting throughout the book
\definecolor{blue}{RGB}{69,74,159} %Redifines the default for the color blue

%----------------------------------------------------------------------------------------
% Font Settings
\usepackage{avant} % Use the Avantgarde font for headings
%\usepackage{times} % Use the Times font for headings
\usepackage{mathptmx} % Use the Adobe Times Roman as the default text font together with math symbols from the Sym­bol, Chancery and Com­puter Modern fonts
\usepackage{microtype} % Slightly tweak font spacing for aesthetics
\usepackage[utf8]{inputenc} % Required for including letters with accents
\usepackage[T1]{fontenc} % Use 8-bit encoding that has 256 glyphs
\usepackage{csquotes}
\usepackage{endnotes}
\usepackage{setspace}
\onehalfspacing

\usepackage{todonotes}

%----------------------------------------------------------------------------------------
% Bibliography
\usepackage[style=authoryear,
            %sorting=nyt,
            %sortcites=true,
            autopunct=true,
            %babel=hyphen,
            %hyperref=true,
            %abbreviate=false,
            backref=true,
            %refsection=chapter,
            backend=bibtex]
{biblatex}
\addbibresource{../../tese.bib} % BibTeX bibliography file
\defbibheading{bibempty}{}
%bibtex sspanel1-blx

%Another option with natbib which require a different method to do a bibliography section per chapter
%\usepackage[square,numbers,sectionbib]{natbib}
%\usepackage{chapterbib}
%\usepackage{authblk}
%\bibliographystyle{evolution}

%----------------------------------------------------------------------------------------
% Index
\usepackage{calc} % For simpler calculation - used for spacing the index letter headings correctly
\usepackage{makeidx} % Required to make an index
\makeindex % Tells LaTeX to create the files required for indexing

%----------------------------------------------------------------------------------------
%----------------------------------------------------------------------------------------
%	VARIOUS REQUIRED PACKAGES
%----------------------------------------------------------------------------------------

\usepackage{titlesec} % Allows customization of titles
\usepackage{graphicx} % Required for including pictures
\graphicspath{{Images/}} % Specifies the directory where pictures are stored
\usepackage{lipsum} % Inserts dummy text
\usepackage{tikz} % Required for drawing custom shapes
\usepackage[brazil]{babel} % Specify language/hyphenation
\usepackage{enumitem} % Customize lists
\setlist{nolistsep} % Reduce spacing between bullet points and numbered lists
\usepackage{booktabs} % Required for nicer horizontal rules in tables
\usepackage{eso-pic} % Required for specifying an image background in the title page
\usepackage[rightcaption]{sidecap} % Required for placing caption on the side of float 

%----------------------------------------------------------------------------------------
%	MAIN TABLE OF CONTENTS
%----------------------------------------------------------------------------------------

\usepackage{titletoc} % Required for manipulating the table of contents

\addtocontents{toc}{\protect\thispagestyle{empty}}
\contentsmargin{0cm} % Removes the default margin
% Chapter text styling
\titlecontents{chapter}[1.25cm] % Indentation
{\addvspace{15pt}\large\sffamily\bfseries} % Spacing and font options for chapters
{\color{color!60}\contentslabel[\Large\thecontentslabel]{1.25cm}\color{color}} % Chapter number
{}  
{\color{color!60}\normalsize\sffamily\bfseries\;\titlerule*[.5pc]{.}\;\thecontentspage} % Page number
% Section text styling
\titlecontents{section}[1.25cm] % Indentation
{\addvspace{5pt}\sffamily\bfseries} % Spacing and font options for sections
{\contentslabel[\thecontentslabel]{1.25cm}} % Section number
{}
{\sffamily\hfill\color{black}\thecontentspage} % Page number
[]
% Subsection text styling
\titlecontents{subsection}[1.25cm] % Indentation
{\addvspace{1pt}\sffamily\small} % Spacing and font options for subsections
{\contentslabel[\thecontentslabel]{1.25cm}} % Subsection number
{}
{\sffamily\;\titlerule*[.5pc]{.}\;\thecontentspage} % Page number
[] 

%----------------------------------------------------------------------------------------
%	MINI TABLE OF CONTENTS IN CHAPTER HEADS
%----------------------------------------------------------------------------------------

% Section text styling
\titlecontents{lsection}[0em] % Indendating
{\footnotesize\sffamily} % Font settings
{}
{}
{}

% Subsection text styling
\titlecontents{lsubsection}[.5em] % Indentation
{\normalfont\footnotesize\sffamily} % Font settings
{}
{}
{}
 
%----------------------------------------------------------------------------------------
%	PAGE HEADERS
%----------------------------------------------------------------------------------------

\usepackage{fancyhdr} % Required for header and footer configuration

\pagestyle{fancy}
\pagenumbering{roman} 
\renewcommand{\chaptermark}[1]{\markboth{\sffamily\normalsize\bfseries #1}{}} % Chapter text font settings
\renewcommand{\sectionmark}[1]{\markright{\sffamily\normalsize\thesection\hspace{5pt}#1}{}} % Section text font settings


%-----------------------------------------------------------------------------------------
%Two-sided printing
%\fancyhf{} \fancyhead[LE,RO]{\sffamily\normalsize\thepage} % Font setting for the page number in the header
%\fancyhead[LO]{\rightmark} % Print the nearest section name on the left side of odd pages
%\fancyhead[RE]{\leftmark} % Print the current chapter name on the right side of even pages
%------------------------------------------------------------------------------------------

%-----------------------------------------------------------------------------------------
%One-sided printing
\fancyhf{} \fancyhead[RE,RO]{\sffamily\normalsize\thepage} % Font setting and placement for the page number in the header
\fancyhead[LO]{\rightmark} % Print the nearest section name on the left side of odd pages
\fancyhead[LE]{\leftmark} % Print the current chapter name on the right side of even pages
%------------------------------------------------------------------------------------------

\renewcommand{\headrulewidth}{0.5pt} % Width of the rule under the header
\addtolength{\headheight}{2.5pt} % Increase the spacing around the header slightly
\renewcommand{\footrulewidth}{0pt} % Removes the rule in the footer
\fancypagestyle{plain}{\fancyfoot[CO,CE]{\thepage}
                      \renewcommand{\headrulewidth}{0pt}} % Style for when a plain pagestyle is specified


\makeatletter
% Removes the header from odd empty pages at the end of chapters
%\renewcommand{\cleardoublepage}{
%\clearpage\ifodd\c@page\else
%\hbox{}
%\vspace*{\fill}
%\thispagestyle{empty}
%\newpage
%\fi}

%----------------------------------------------------------------------------------------
%	SET UP POSITION OF TEXT IN TITLE PAGE
%----------------------------------------------------------------------------------------
\def\parsecomma#1,#2\endparsecomma{\def\page@x{#1}\def\page@y{#2}}
\tikzdeclarecoordinatesystem{page}{
    \parsecomma#1\endparsecomma
    \pgfpointanchor{current page}{north east}
    % Save the upper right corner
    \pgf@xc=\pgf@x%
    \pgf@yc=\pgf@y%
    % save the lower left corner
    \pgfpointanchor{current page}{south west}
    \pgf@xb=\pgf@x%
    \pgf@yb=\pgf@y%
    % Transform to the correct placement
    \pgfmathparse{(\pgf@xc-\pgf@xb)/2.*\page@x+(\pgf@xc+\pgf@xb)/2.}
    \expandafter\pgf@x\expandafter=\pgfmathresult pt
    \pgfmathparse{(\pgf@yc-\pgf@yb)/2.*\page@y+(\pgf@yc+\pgf@yb)/2.}
    \expandafter\pgf@y\expandafter=\pgfmathresult pt
}

%----------------------------------------------------------------------------------------
%	MODIFIED ENDNOTE AND FOOTNOTES COMMANDS
%----------------------------------------------------------------------------------------
\renewcommand{\theendnote}{\alph{endnote}}

\newcommand\footnoteref[1]{\protected@xdef\@thefnmark{\ref{#1}}\@footnotemark}

%Use the following command to add a endnote to a footnote without showing the number.

\newcommand\numberlessnote[1]{%
  \begingroup
  \renewcommand\thefootnote{}\footnote{#1}%
  \addtocounter{footnote}{-1}%
  \endgroup
}

%\newcommand\insidenote[1]{%
%  \begingroup
%  \addtocounter{endnote}{+1}%
%  \renewcommand\theendnotes{}\endnotetext{#1}%
%  \endgroup
%}

%Use the following command if you are using insidenotes in our text and one to keep which
%note number unique
%Example:
% Some text \footnote{See page \ref{endnotes} \insidenote{Actual note}}
% The label endnote should then be placed in the page with the notes.

%\newcommand\onlyfootnote[1]{%
%  \begingroup
%  \addtocounter{endnote}{+1}%
%  \renewcommand\theendnotes{}\footnote{#1}%
%  \endgroup
%}

%\renewcommand*{\thefootnote}{\fnsymbol{footnote}}

%----------------------------------------------------------------------------------------
%	MODIFIED TABLES
%----------------------------------------------------------------------------------------

\usepackage{tabularx}
\newcolumntype{C}[1]{>{\hsize=#1\hsize\centering\arraybackslash}X}%
\newcolumntype{L}[1]{>{\hsize=#1\hsize\raggedright\arraybackslash}X}%


\def\hlinewd#1{%
\noalign{\ifnum0=`}\fi\hrule \@height #1 %
\futurelet\reserved@a\@xhline}


\newcommand{\abstractseparator}{
\textcolor{color}{\rule[0.5ex]{1\columnwidth}{2pt}}
}

%----------------------------------------------------------------------------------------
%	THEOREM STYLES
%----------------------------------------------------------------------------------------


\usepackage{amsmath,amsfonts,amssymb,amsthm} % For including math equations, theorems, symbols, etc

\newcommand{\intoo}[2]{\mathopen{]}#1\,;#2\mathclose{[}}
\newcommand{\ud}{\mathop{\mathrm{{}d}}\mathopen{}}
\newcommand{\intff}[2]{\mathopen{[}#1\,;#2\mathclose{]}}
\newtheorem{notation}{Notation}[chapter]

%%%%%%%%%%%%%%%%%%%%%%%%%%%%%%%%%%%%%%%%%%%%%%%%%%%%%%%%%%%%%%%%%%%%%%%%%%%
%%%%%%%%%%%%%%%%%%%% dedicated to boxed/framed environements %%%%%%%%%%%%%%
%%%%%%%%%%%%%%%%%%%%%%%%%%%%%%%%%%%%%%%%%%%%%%%%%%%%%%%%%%%%%%%%%%%%%%%%%%%

\newtheoremstyle{ocrenumbox}% % Theorem style name
{0pt}% Space above
{0pt}% Space below
{\normalfont}% % Body font
{}% Indent amount
{\small\bf\sffamily\color{color}}% % Theorem head font
{\;}% Punctuation after theorem head
{0.25em}% Space after theorem head
{\small\sffamily\color{color}\thmname{#1}\nobreakspace\thmnumber{\@ifnotempty{#1}{}\@upn{#2}}% Theorem text (e.g. Theorem 2.1)
\thmnote{\nobreakspace\the\thm@notefont\sffamily\bfseries\color{black}---\nobreakspace#3.}} % Optional theorem note
\renewcommand{\qedsymbol}{$\blacksquare$}% Optional qed square

\newtheoremstyle{blacknumex}% Theorem style name
{5pt}% Space above
{5pt}% Space below
{\normalfont}% Body font
{} % Indent amount
{\small\bf\sffamily}% Theorem head font
{\;}% Punctuation after theorem head
{0.25em}% Space after theorem head
{\small\sffamily{\tiny\ensuremath{\blacksquare}}\nobreakspace\thmname{#1}\nobreakspace\thmnumber{\@ifnotempty{#1}{}\@upn{#2}}% Theorem text (e.g. Theorem 2.1)
\thmnote{\nobreakspace\the\thm@notefont\sffamily\bfseries---\nobreakspace#3.}}% Optional theorem note

\newtheoremstyle{blacknumbox} % Theorem style name
{0pt}% Space above
{0pt}% Space below
{\normalfont}% Body font
{}% Indent amount
{\small\bf\sffamily}% Theorem head font
{\;}% Punctuation after theorem head
{0.25em}% Space after theorem head
{\small\sffamily\thmname{#1}\nobreakspace\thmnumber{\@ifnotempty{#1}{}\@upn{#2}}% Theorem text (e.g. Theorem 2.1)
\thmnote{\nobreakspace\the\thm@notefont\sffamily\bfseries---\nobreakspace#3.}}% Optional theorem note


\newtheoremstyle{coloredbox}% Box style name
{0pt}% Space above
{0pt}% Space below
{\normalfont}% % Body font
{}% Indent amount
{\small\bf\sffamily\color{color}}% % Theorem head font
{\;}% Punctuation after theorem head
{\newline}% Space after theorem head
{\small\sffamily\color{color}\thmname{#1}\nobreakspace\thmnumber{\@ifnotempty{#1}{}\@upn{#2}}% Theorem text (e.g. Theorem 2.1)
\thmnote{\nobreakspace\the\thm@notefont\sffamily\bfseries\color{black}---\nobreakspace#3.}} % Optional theorem note
\renewcommand{\qedsymbol}{$\blacksquare$}% Optional qed square

%%%%%%%%%%%%%%%%%%%%%%%%%%%%%%%%%%%%%%%%%%%%%%%%%%%%%%%%%%%%%%%%%%%%%%%%%%%
%%%%%%%%%%%%% dedicated to non-boxed/non-framed environements %%%%%%%%%%%%%
%%%%%%%%%%%%%%%%%%%%%%%%%%%%%%%%%%%%%%%%%%%%%%%%%%%%%%%%%%%%%%%%%%%%%%%%%%%
\newtheoremstyle{ocrenum}% % Theorem style name
{5pt}% Space above
{5pt}% Space below
{\normalfont}% % Body font
{}% Indent amount
{\small\bf\sffamily\color{color}}% % Theorem head font
{\;}% Punctuation after theorem head
{0.25em}% Space after theorem head
{\small\sffamily\color{color}\thmname{#1}\nobreakspace\thmnumber{\@ifnotempty{#1}{}\@upn{#2}}% Theorem text (e.g. Theorem 2.1)
\thmnote{\nobreakspace\the\thm@notefont\sffamily\bfseries\color{black}---\nobreakspace#3.}} % Optional theorem note
\renewcommand{\qedsymbol}{$\blacksquare$}% Optional qed square
\makeatother

% Defines the theorem text style for each type of theorem to one of the three styles above
\newcounter{dummy} 
\numberwithin{dummy}{section}
\theoremstyle{ocrenumbox}
\newtheorem{theoremeT}[dummy]{Theorem}
\newtheorem{problem}{Problem}[chapter]
\newtheorem{exerciseT}{Exercise}[chapter]
\theoremstyle{blacknumex}
\newtheorem{exampleT}{Example}[chapter]
\theoremstyle{blacknumbox}
\newtheorem{vocabulary}{Vocabulary}[chapter]
\newtheorem{definitionT}{Definition}[section]
\newtheorem{corollaryT}[dummy]{Corollary}
\theoremstyle{ocrenum}
\newtheorem{proposition}[dummy]{Proposition}

\theoremstyle{coloredbox}
\newtheorem{quadroT}{Quadro}[chapter]
%\newtheorem{quadroT}{Box}[chapter] %english

%----------------------------------------------------------------------------------------
%	DEFINITION OF COLORED BOXES
%----------------------------------------------------------------------------------------

\RequirePackage[framemethod=default]{mdframed} % Required for creating the theorem, definition, exercise and corollary boxes

% Theorem box
\newmdenv[skipabove=7pt,
skipbelow=7pt,
backgroundcolor=black!5,
linecolor=color,
innerleftmargin=5pt,
innerrightmargin=5pt,
innertopmargin=5pt,
leftmargin=0cm,
rightmargin=0cm,
innerbottommargin=5pt]{tBox}

% Exercise box	  
\newmdenv[skipabove=7pt,
skipbelow=7pt,
rightline=false,
leftline=true,
topline=false,
bottomline=false,
backgroundcolor=color!10,
linecolor=color,
innerleftmargin=5pt,
innerrightmargin=5pt,
innertopmargin=5pt,
innerbottommargin=5pt,
leftmargin=0cm,
rightmargin=0cm,
linewidth=4pt]{eBox}

% Quadro	  
\newmdenv[skipabove=7pt,
skipbelow=7pt,
rightline=false,
leftline=true,
topline=false,
bottomline=false,
backgroundcolor=color!30,
linecolor=color,
innerleftmargin=5pt,
innerrightmargin=5pt,
innertopmargin=5pt,
innerbottommargin=5pt,
leftmargin=0cm,
rightmargin=0cm,
linewidth=4pt]{iBox}

% Quadro branco  
\newmdenv[skipabove=7pt,
skipbelow=7pt,
rightline=true,
leftline=true,
topline=true,
bottomline=true,
backgroundcolor=white,
linecolor=color,
innerleftmargin=15pt,
innerrightmargin=15pt,
innertopmargin=5pt,
innerbottommargin=5pt,
leftmargin=0cm,
rightmargin=0cm,
linewidth=2pt]{clearBox}

% Definition box
\newmdenv[skipabove=7pt,
skipbelow=7pt,
rightline=false,
leftline=true,
topline=false,
bottomline=false,
linecolor=color,
innerleftmargin=5pt,
innerrightmargin=5pt,
innertopmargin=0pt,
leftmargin=0cm,
rightmargin=0cm,
linewidth=4pt,
innerbottommargin=0pt]{dBox}	

% Corollary box
\newmdenv[skipabove=7pt,
skipbelow=7pt,
rightline=false,
leftline=true,
topline=false,
bottomline=false,
linecolor=gray,
backgroundcolor=black!5,
innerleftmargin=5pt,
innerrightmargin=5pt,
innertopmargin=5pt,
leftmargin=0cm,
rightmargin=0cm,
linewidth=4pt,
innerbottommargin=5pt]{cBox}				  
		  

% Creates an environment for each type of theorem and assigns it a theorem text style from the "Theorem Styles" section above and a colored box from above
\newenvironment{theorem}{\begin{tBox}\begin{theoremeT}}{\end{theoremeT}\end{tBox}}
\newenvironment{exercise}{\begin{eBox}\begin{exerciseT}}{\hfill{\color{color}\tiny\ensuremath{\blacksquare}}\end{exerciseT}\end{eBox}}				  
\newenvironment{definition}{\begin{dBox}\begin{definitionT}}{\end{definitionT}\end{dBox}}	
\newenvironment{example}{\begin{exampleT}}{\hfill{\tiny\ensuremath{\blacksquare}}\end{exampleT}}		
\newenvironment{corollary}{\begin{cBox}\begin{corollaryT}}{\end{corollaryT}\end{cBox}}

\newenvironment{infobox}
  {\begin{iBox}\begin{quadroT}}
    {\hfill{\color{color}\tiny\ensuremath{\blacksquare}}
   \end{quadroT}\end{iBox}}  
   
\newenvironment{clearinfobox}
  {\begin{clearBox}\begin{quadroT}}
    {\hfill{\color{color}\tiny\ensuremath{\blacksquare}}
   \end{quadroT}\end{clearBox}}    


%----------------------------------------------------------------------------------------
%	REMARK ENVIRONMENT
%----------------------------------------------------------------------------------------

\newenvironment{remark}{\par\vskip10pt\small % Vertical white space above the remark and smaller font size
\begin{list}{}{
\leftmargin=35pt % Indentation on the left
\rightmargin=25pt}\item\ignorespaces % Indentation on the right
\makebox[-2.5pt]{\begin{tikzpicture}[overlay]
\node[draw=color!60,line width=1pt,circle,fill=color!25,font=\sffamily\bfseries,inner sep=2pt,outer sep=0pt] at (-15pt,0pt){\textcolor{color}{N}};\end{tikzpicture}} % Orange R in a circle
\advance\baselineskip -1pt}{\end{list}\vskip5pt} % Tighter line spacing and white space after remark

%----------------------------------------------------------------------------------------
%	SECTION NUMBERING IN THE MARGIN
%----------------------------------------------------------------------------------------

\makeatletter
\renewcommand{\@seccntformat}[1]{\llap{\textcolor{color}{\csname the#1\endcsname}\hspace{1em}}}                    
\renewcommand{\section}{\@startsection{section}{1}{\z@}
{-4ex \@plus -1ex \@minus -.4ex}
{1ex \@plus.2ex }
{\normalfont\large\sffamily\bfseries}}
\renewcommand{\subsection}{\@startsection {subsection}{2}{\z@}
{-3ex \@plus -0.1ex \@minus -.4ex}
{0.5ex \@plus.2ex }
{\normalfont\sffamily\bfseries}}
\renewcommand{\subsubsection}{\@startsection {subsubsection}{3}{\z@}
{-2ex \@plus -0.1ex \@minus -.2ex}
{0.2ex \@plus.2ex }
{\normalfont\small\sffamily\bfseries}}                        
\renewcommand\paragraph{\@startsection{paragraph}{4}{\z@}
{-2ex \@plus-.2ex \@minus .2ex}
{0.1ex}
{\normalfont\small\sffamily\bfseries}}

%----------------------------------------------------------------------------------------
%	CHAPTER HEADINGS
%----------------------------------------------------------------------------------------

\newcommand{\thechapterimage}{}
\newcommand{\chapterimage}[1]{\renewcommand{\thechapterimage}{#1}}
\def\thechapter{\arabic{chapter}}
\def\@makechapterhead#1{
\thispagestyle{empty}
{\centering \normalfont\sffamily
\ifnum \c@secnumdepth >\m@ne
\if@mainmatter
\startcontents
\begin{tikzpicture}[remember picture,overlay]
\node at (current page.north west)
{\begin{tikzpicture}[remember picture,overlay]

\node[anchor=north west,inner sep=0pt] at (0,0) {\includegraphics[width=\paperwidth]{\thechapterimage}};

%Commenting the 3 lines below removes the small contents box in the chapter heading
\draw[fill=white,opacity=1] (1cm, 0) rectangle (8cm,-7cm);
\node[anchor=north west] at (1cm,.25cm) {\parbox[t][8cm][t]{6.5cm}{\huge\bfseries\flushleft \printcontents{l}{1}{\setcounter{tocdepth}{2}}}};
%\node[anchor=north west] at (1cm,.25cm) {\parbox[t][8cm][t]{6.5cm}{\huge\bfseries\flushleft \printcontents{l}{1}{\setcounter{tocdepth}{1}}}};
\draw[anchor=west] (5cm,-9cm) node [rounded corners=25pt,fill=white,fill opacity=1,text opacity=1,draw=color,draw opacity=1,line width=2pt,inner sep=15pt]{\huge\sffamily\bfseries\textcolor{black}{\thechapter\ ---\ #1\vphantom{plPQq}\makebox[22cm]{}}};

\end{tikzpicture}};
\end{tikzpicture}}\par\vspace*{230\p@}
\fi
\fi
}
\def\@makeschapterhead#1{
\thispagestyle{empty}
{\centering \normalfont\sffamily
\ifnum \c@secnumdepth >\m@ne
\if@mainmatter
%\startcontents
\begin{tikzpicture}[remember picture,overlay]
\node at (current page.north west)
{\begin{tikzpicture}[remember picture,overlay]
%\node[anchor=north west,inner sep=0pt] at (0,0) {\includegraphics[width=\paperwidth]{\thechapterimage}};
%Commenting the 3 lines below removes the small contents box in the chapter heading
%\draw[fill=white,opacity=1] (1cm,0) rectangle (8cm,-7cm);
%\node[anchor=north west] at (1cm,.25cm) {\parbox[t][8cm][t]{6.5cm}{\huge\bfseries\flushleft \printcontents{l}{1}{\setcounter{tocdepth}{2}}}};
\draw[anchor=west] (-0.2cm,-2.5cm) node [fill=white,fill opacity=1,text opacity=1,draw=color,draw opacity=1,line width=2pt,inner sep=15pt, inner xsep=60pt]{\huge\sffamily\bfseries\textcolor{black}{  \ #1\vphantom{plPQq}\makebox[22cm]{}}};
\end{tikzpicture}};
\end{tikzpicture}}\par\vspace*{50\p@}
\fi
\fi
}
\makeatother % Insert the commands.tex file which contains the majority of the structure behind the template

%----------------------------------------------------------------------------------------

\begin{document}
\graphicspath{ {../../Images/}}
\fi


\chapterimage{Headers/chapter_header.pdf}
\chapter[Research article]{Short title
\label{chap:article}}
%\begin{refsection}
%------------------------------------------------
%\listoftodos

{\Large\sffamily\textbf {Título}}\\
\vspace{0.6cm}
\begin{tikzpicture}[remember picture,overlay]
     \node[white,draw,text=black,text width=11.5cm] at (page cs:0.15,0.1) 
     {\it\footnotesize O conteúdo deste capítulo foi publicado no artigo:
      \textbf{Autores}. ANO. 
      “Título do artigo”. Revista, volume, páginas. doi.      
      };
\end{tikzpicture}
\vspace{0.5cm}
\newline
\abstractseparator
{\large\sffamily\textbf {Resumo}}\\
{Resumo de capítulo/artigo.}\\
\abstractseparator

\section{Specific introduction topic}

\lipsum[1-6]

\section{ETC.}



%\fi
%------------------------------------------------
%\printbibliography[heading=subbibliography]
%\end{refsection}

\iffalse
\begingroup
\let\enotesize\normalsize
\renewcommand\notesname{Notas finais}
\renewcommand\enoteformat{\noindent[\theenmark]\hspace{0.06in}}
\label{notas}
\theendnotes
\endgroup

\end{document}
\fi

%----------------------------------------------------------------------------------------
%	CAPÍTULO 3
%----------------------------------------------------------------------------------------
\iffalse
\documentclass[11pt,fleqn]{book} % Default font size and left-justified equations
\usepackage[top=3cm,bottom=3cm,left=3.2cm,right=3.2cm,headsep=10pt,a4paper]{geometry} % Page margins
\usepackage{xcolor} % Required for specifying colors by name
\definecolor{color}{RGB}{128,0,0} % Define the color used for highlighting throughout the book
\definecolor{blue}{RGB}{69,74,159} %Redifines the default for the color blue

%----------------------------------------------------------------------------------------
% Font Settings
\usepackage{avant} % Use the Avantgarde font for headings
%\usepackage{times} % Use the Times font for headings
\usepackage{mathptmx} % Use the Adobe Times Roman as the default text font together with math symbols from the Sym­bol, Chancery and Com­puter Modern fonts
\usepackage{microtype} % Slightly tweak font spacing for aesthetics
\usepackage[utf8]{inputenc} % Required for including letters with accents
\usepackage[T1]{fontenc} % Use 8-bit encoding that has 256 glyphs
\usepackage{csquotes}
\usepackage{endnotes}
\usepackage{setspace}
\onehalfspacing

\usepackage{todonotes}

%----------------------------------------------------------------------------------------
% Bibliography
\usepackage[style=authoryear,
            %sorting=nyt,
            %sortcites=true,
            autopunct=true,
            %babel=hyphen,
            %hyperref=true,
            %abbreviate=false,
            backref=true,
            %refsection=chapter,
            backend=bibtex]
{biblatex}
\addbibresource{../../tese.bib} % BibTeX bibliography file
\defbibheading{bibempty}{}
%bibtex sspanel1-blx

%Another option with natbib which require a different method to do a bibliography section per chapter
%\usepackage[square,numbers,sectionbib]{natbib}
%\usepackage{chapterbib}
%\usepackage{authblk}
%\bibliographystyle{evolution}

%----------------------------------------------------------------------------------------
% Index
\usepackage{calc} % For simpler calculation - used for spacing the index letter headings correctly
\usepackage{makeidx} % Required to make an index
\makeindex % Tells LaTeX to create the files required for indexing

%----------------------------------------------------------------------------------------
%----------------------------------------------------------------------------------------
%	VARIOUS REQUIRED PACKAGES
%----------------------------------------------------------------------------------------

\usepackage{titlesec} % Allows customization of titles
\usepackage{graphicx} % Required for including pictures
\graphicspath{{Images/}} % Specifies the directory where pictures are stored
\usepackage{lipsum} % Inserts dummy text
\usepackage{tikz} % Required for drawing custom shapes
\usepackage[brazil]{babel} % Specify language/hyphenation
\usepackage{enumitem} % Customize lists
\setlist{nolistsep} % Reduce spacing between bullet points and numbered lists
\usepackage{booktabs} % Required for nicer horizontal rules in tables
\usepackage{eso-pic} % Required for specifying an image background in the title page
\usepackage[rightcaption]{sidecap} % Required for placing caption on the side of float 

%----------------------------------------------------------------------------------------
%	MAIN TABLE OF CONTENTS
%----------------------------------------------------------------------------------------

\usepackage{titletoc} % Required for manipulating the table of contents

\addtocontents{toc}{\protect\thispagestyle{empty}}
\contentsmargin{0cm} % Removes the default margin
% Chapter text styling
\titlecontents{chapter}[1.25cm] % Indentation
{\addvspace{15pt}\large\sffamily\bfseries} % Spacing and font options for chapters
{\color{color!60}\contentslabel[\Large\thecontentslabel]{1.25cm}\color{color}} % Chapter number
{}  
{\color{color!60}\normalsize\sffamily\bfseries\;\titlerule*[.5pc]{.}\;\thecontentspage} % Page number
% Section text styling
\titlecontents{section}[1.25cm] % Indentation
{\addvspace{5pt}\sffamily\bfseries} % Spacing and font options for sections
{\contentslabel[\thecontentslabel]{1.25cm}} % Section number
{}
{\sffamily\hfill\color{black}\thecontentspage} % Page number
[]
% Subsection text styling
\titlecontents{subsection}[1.25cm] % Indentation
{\addvspace{1pt}\sffamily\small} % Spacing and font options for subsections
{\contentslabel[\thecontentslabel]{1.25cm}} % Subsection number
{}
{\sffamily\;\titlerule*[.5pc]{.}\;\thecontentspage} % Page number
[] 

%----------------------------------------------------------------------------------------
%	MINI TABLE OF CONTENTS IN CHAPTER HEADS
%----------------------------------------------------------------------------------------

% Section text styling
\titlecontents{lsection}[0em] % Indendating
{\footnotesize\sffamily} % Font settings
{}
{}
{}

% Subsection text styling
\titlecontents{lsubsection}[.5em] % Indentation
{\normalfont\footnotesize\sffamily} % Font settings
{}
{}
{}
 
%----------------------------------------------------------------------------------------
%	PAGE HEADERS
%----------------------------------------------------------------------------------------

\usepackage{fancyhdr} % Required for header and footer configuration

\pagestyle{fancy}
\pagenumbering{roman} 
\renewcommand{\chaptermark}[1]{\markboth{\sffamily\normalsize\bfseries #1}{}} % Chapter text font settings
\renewcommand{\sectionmark}[1]{\markright{\sffamily\normalsize\thesection\hspace{5pt}#1}{}} % Section text font settings


%-----------------------------------------------------------------------------------------
%Two-sided printing
%\fancyhf{} \fancyhead[LE,RO]{\sffamily\normalsize\thepage} % Font setting for the page number in the header
%\fancyhead[LO]{\rightmark} % Print the nearest section name on the left side of odd pages
%\fancyhead[RE]{\leftmark} % Print the current chapter name on the right side of even pages
%------------------------------------------------------------------------------------------

%-----------------------------------------------------------------------------------------
%One-sided printing
\fancyhf{} \fancyhead[RE,RO]{\sffamily\normalsize\thepage} % Font setting and placement for the page number in the header
\fancyhead[LO]{\rightmark} % Print the nearest section name on the left side of odd pages
\fancyhead[LE]{\leftmark} % Print the current chapter name on the right side of even pages
%------------------------------------------------------------------------------------------

\renewcommand{\headrulewidth}{0.5pt} % Width of the rule under the header
\addtolength{\headheight}{2.5pt} % Increase the spacing around the header slightly
\renewcommand{\footrulewidth}{0pt} % Removes the rule in the footer
\fancypagestyle{plain}{\fancyfoot[CO,CE]{\thepage}
                      \renewcommand{\headrulewidth}{0pt}} % Style for when a plain pagestyle is specified


\makeatletter
% Removes the header from odd empty pages at the end of chapters
%\renewcommand{\cleardoublepage}{
%\clearpage\ifodd\c@page\else
%\hbox{}
%\vspace*{\fill}
%\thispagestyle{empty}
%\newpage
%\fi}

%----------------------------------------------------------------------------------------
%	SET UP POSITION OF TEXT IN TITLE PAGE
%----------------------------------------------------------------------------------------
\def\parsecomma#1,#2\endparsecomma{\def\page@x{#1}\def\page@y{#2}}
\tikzdeclarecoordinatesystem{page}{
    \parsecomma#1\endparsecomma
    \pgfpointanchor{current page}{north east}
    % Save the upper right corner
    \pgf@xc=\pgf@x%
    \pgf@yc=\pgf@y%
    % save the lower left corner
    \pgfpointanchor{current page}{south west}
    \pgf@xb=\pgf@x%
    \pgf@yb=\pgf@y%
    % Transform to the correct placement
    \pgfmathparse{(\pgf@xc-\pgf@xb)/2.*\page@x+(\pgf@xc+\pgf@xb)/2.}
    \expandafter\pgf@x\expandafter=\pgfmathresult pt
    \pgfmathparse{(\pgf@yc-\pgf@yb)/2.*\page@y+(\pgf@yc+\pgf@yb)/2.}
    \expandafter\pgf@y\expandafter=\pgfmathresult pt
}

%----------------------------------------------------------------------------------------
%	MODIFIED ENDNOTE AND FOOTNOTES COMMANDS
%----------------------------------------------------------------------------------------
\renewcommand{\theendnote}{\alph{endnote}}

\newcommand\footnoteref[1]{\protected@xdef\@thefnmark{\ref{#1}}\@footnotemark}

%Use the following command to add a endnote to a footnote without showing the number.

\newcommand\numberlessnote[1]{%
  \begingroup
  \renewcommand\thefootnote{}\footnote{#1}%
  \addtocounter{footnote}{-1}%
  \endgroup
}

%\newcommand\insidenote[1]{%
%  \begingroup
%  \addtocounter{endnote}{+1}%
%  \renewcommand\theendnotes{}\endnotetext{#1}%
%  \endgroup
%}

%Use the following command if you are using insidenotes in our text and one to keep which
%note number unique
%Example:
% Some text \footnote{See page \ref{endnotes} \insidenote{Actual note}}
% The label endnote should then be placed in the page with the notes.

%\newcommand\onlyfootnote[1]{%
%  \begingroup
%  \addtocounter{endnote}{+1}%
%  \renewcommand\theendnotes{}\footnote{#1}%
%  \endgroup
%}

%\renewcommand*{\thefootnote}{\fnsymbol{footnote}}

%----------------------------------------------------------------------------------------
%	MODIFIED TABLES
%----------------------------------------------------------------------------------------

\usepackage{tabularx}
\newcolumntype{C}[1]{>{\hsize=#1\hsize\centering\arraybackslash}X}%
\newcolumntype{L}[1]{>{\hsize=#1\hsize\raggedright\arraybackslash}X}%


\def\hlinewd#1{%
\noalign{\ifnum0=`}\fi\hrule \@height #1 %
\futurelet\reserved@a\@xhline}


\newcommand{\abstractseparator}{
\textcolor{color}{\rule[0.5ex]{1\columnwidth}{2pt}}
}

%----------------------------------------------------------------------------------------
%	THEOREM STYLES
%----------------------------------------------------------------------------------------


\usepackage{amsmath,amsfonts,amssymb,amsthm} % For including math equations, theorems, symbols, etc

\newcommand{\intoo}[2]{\mathopen{]}#1\,;#2\mathclose{[}}
\newcommand{\ud}{\mathop{\mathrm{{}d}}\mathopen{}}
\newcommand{\intff}[2]{\mathopen{[}#1\,;#2\mathclose{]}}
\newtheorem{notation}{Notation}[chapter]

%%%%%%%%%%%%%%%%%%%%%%%%%%%%%%%%%%%%%%%%%%%%%%%%%%%%%%%%%%%%%%%%%%%%%%%%%%%
%%%%%%%%%%%%%%%%%%%% dedicated to boxed/framed environements %%%%%%%%%%%%%%
%%%%%%%%%%%%%%%%%%%%%%%%%%%%%%%%%%%%%%%%%%%%%%%%%%%%%%%%%%%%%%%%%%%%%%%%%%%

\newtheoremstyle{ocrenumbox}% % Theorem style name
{0pt}% Space above
{0pt}% Space below
{\normalfont}% % Body font
{}% Indent amount
{\small\bf\sffamily\color{color}}% % Theorem head font
{\;}% Punctuation after theorem head
{0.25em}% Space after theorem head
{\small\sffamily\color{color}\thmname{#1}\nobreakspace\thmnumber{\@ifnotempty{#1}{}\@upn{#2}}% Theorem text (e.g. Theorem 2.1)
\thmnote{\nobreakspace\the\thm@notefont\sffamily\bfseries\color{black}---\nobreakspace#3.}} % Optional theorem note
\renewcommand{\qedsymbol}{$\blacksquare$}% Optional qed square

\newtheoremstyle{blacknumex}% Theorem style name
{5pt}% Space above
{5pt}% Space below
{\normalfont}% Body font
{} % Indent amount
{\small\bf\sffamily}% Theorem head font
{\;}% Punctuation after theorem head
{0.25em}% Space after theorem head
{\small\sffamily{\tiny\ensuremath{\blacksquare}}\nobreakspace\thmname{#1}\nobreakspace\thmnumber{\@ifnotempty{#1}{}\@upn{#2}}% Theorem text (e.g. Theorem 2.1)
\thmnote{\nobreakspace\the\thm@notefont\sffamily\bfseries---\nobreakspace#3.}}% Optional theorem note

\newtheoremstyle{blacknumbox} % Theorem style name
{0pt}% Space above
{0pt}% Space below
{\normalfont}% Body font
{}% Indent amount
{\small\bf\sffamily}% Theorem head font
{\;}% Punctuation after theorem head
{0.25em}% Space after theorem head
{\small\sffamily\thmname{#1}\nobreakspace\thmnumber{\@ifnotempty{#1}{}\@upn{#2}}% Theorem text (e.g. Theorem 2.1)
\thmnote{\nobreakspace\the\thm@notefont\sffamily\bfseries---\nobreakspace#3.}}% Optional theorem note


\newtheoremstyle{coloredbox}% Box style name
{0pt}% Space above
{0pt}% Space below
{\normalfont}% % Body font
{}% Indent amount
{\small\bf\sffamily\color{color}}% % Theorem head font
{\;}% Punctuation after theorem head
{\newline}% Space after theorem head
{\small\sffamily\color{color}\thmname{#1}\nobreakspace\thmnumber{\@ifnotempty{#1}{}\@upn{#2}}% Theorem text (e.g. Theorem 2.1)
\thmnote{\nobreakspace\the\thm@notefont\sffamily\bfseries\color{black}---\nobreakspace#3.}} % Optional theorem note
\renewcommand{\qedsymbol}{$\blacksquare$}% Optional qed square

%%%%%%%%%%%%%%%%%%%%%%%%%%%%%%%%%%%%%%%%%%%%%%%%%%%%%%%%%%%%%%%%%%%%%%%%%%%
%%%%%%%%%%%%% dedicated to non-boxed/non-framed environements %%%%%%%%%%%%%
%%%%%%%%%%%%%%%%%%%%%%%%%%%%%%%%%%%%%%%%%%%%%%%%%%%%%%%%%%%%%%%%%%%%%%%%%%%
\newtheoremstyle{ocrenum}% % Theorem style name
{5pt}% Space above
{5pt}% Space below
{\normalfont}% % Body font
{}% Indent amount
{\small\bf\sffamily\color{color}}% % Theorem head font
{\;}% Punctuation after theorem head
{0.25em}% Space after theorem head
{\small\sffamily\color{color}\thmname{#1}\nobreakspace\thmnumber{\@ifnotempty{#1}{}\@upn{#2}}% Theorem text (e.g. Theorem 2.1)
\thmnote{\nobreakspace\the\thm@notefont\sffamily\bfseries\color{black}---\nobreakspace#3.}} % Optional theorem note
\renewcommand{\qedsymbol}{$\blacksquare$}% Optional qed square
\makeatother

% Defines the theorem text style for each type of theorem to one of the three styles above
\newcounter{dummy} 
\numberwithin{dummy}{section}
\theoremstyle{ocrenumbox}
\newtheorem{theoremeT}[dummy]{Theorem}
\newtheorem{problem}{Problem}[chapter]
\newtheorem{exerciseT}{Exercise}[chapter]
\theoremstyle{blacknumex}
\newtheorem{exampleT}{Example}[chapter]
\theoremstyle{blacknumbox}
\newtheorem{vocabulary}{Vocabulary}[chapter]
\newtheorem{definitionT}{Definition}[section]
\newtheorem{corollaryT}[dummy]{Corollary}
\theoremstyle{ocrenum}
\newtheorem{proposition}[dummy]{Proposition}

\theoremstyle{coloredbox}
\newtheorem{quadroT}{Quadro}[chapter]
%\newtheorem{quadroT}{Box}[chapter] %english

%----------------------------------------------------------------------------------------
%	DEFINITION OF COLORED BOXES
%----------------------------------------------------------------------------------------

\RequirePackage[framemethod=default]{mdframed} % Required for creating the theorem, definition, exercise and corollary boxes

% Theorem box
\newmdenv[skipabove=7pt,
skipbelow=7pt,
backgroundcolor=black!5,
linecolor=color,
innerleftmargin=5pt,
innerrightmargin=5pt,
innertopmargin=5pt,
leftmargin=0cm,
rightmargin=0cm,
innerbottommargin=5pt]{tBox}

% Exercise box	  
\newmdenv[skipabove=7pt,
skipbelow=7pt,
rightline=false,
leftline=true,
topline=false,
bottomline=false,
backgroundcolor=color!10,
linecolor=color,
innerleftmargin=5pt,
innerrightmargin=5pt,
innertopmargin=5pt,
innerbottommargin=5pt,
leftmargin=0cm,
rightmargin=0cm,
linewidth=4pt]{eBox}

% Quadro	  
\newmdenv[skipabove=7pt,
skipbelow=7pt,
rightline=false,
leftline=true,
topline=false,
bottomline=false,
backgroundcolor=color!30,
linecolor=color,
innerleftmargin=5pt,
innerrightmargin=5pt,
innertopmargin=5pt,
innerbottommargin=5pt,
leftmargin=0cm,
rightmargin=0cm,
linewidth=4pt]{iBox}

% Quadro branco  
\newmdenv[skipabove=7pt,
skipbelow=7pt,
rightline=true,
leftline=true,
topline=true,
bottomline=true,
backgroundcolor=white,
linecolor=color,
innerleftmargin=15pt,
innerrightmargin=15pt,
innertopmargin=5pt,
innerbottommargin=5pt,
leftmargin=0cm,
rightmargin=0cm,
linewidth=2pt]{clearBox}

% Definition box
\newmdenv[skipabove=7pt,
skipbelow=7pt,
rightline=false,
leftline=true,
topline=false,
bottomline=false,
linecolor=color,
innerleftmargin=5pt,
innerrightmargin=5pt,
innertopmargin=0pt,
leftmargin=0cm,
rightmargin=0cm,
linewidth=4pt,
innerbottommargin=0pt]{dBox}	

% Corollary box
\newmdenv[skipabove=7pt,
skipbelow=7pt,
rightline=false,
leftline=true,
topline=false,
bottomline=false,
linecolor=gray,
backgroundcolor=black!5,
innerleftmargin=5pt,
innerrightmargin=5pt,
innertopmargin=5pt,
leftmargin=0cm,
rightmargin=0cm,
linewidth=4pt,
innerbottommargin=5pt]{cBox}				  
		  

% Creates an environment for each type of theorem and assigns it a theorem text style from the "Theorem Styles" section above and a colored box from above
\newenvironment{theorem}{\begin{tBox}\begin{theoremeT}}{\end{theoremeT}\end{tBox}}
\newenvironment{exercise}{\begin{eBox}\begin{exerciseT}}{\hfill{\color{color}\tiny\ensuremath{\blacksquare}}\end{exerciseT}\end{eBox}}				  
\newenvironment{definition}{\begin{dBox}\begin{definitionT}}{\end{definitionT}\end{dBox}}	
\newenvironment{example}{\begin{exampleT}}{\hfill{\tiny\ensuremath{\blacksquare}}\end{exampleT}}		
\newenvironment{corollary}{\begin{cBox}\begin{corollaryT}}{\end{corollaryT}\end{cBox}}

\newenvironment{infobox}
  {\begin{iBox}\begin{quadroT}}
    {\hfill{\color{color}\tiny\ensuremath{\blacksquare}}
   \end{quadroT}\end{iBox}}  
   
\newenvironment{clearinfobox}
  {\begin{clearBox}\begin{quadroT}}
    {\hfill{\color{color}\tiny\ensuremath{\blacksquare}}
   \end{quadroT}\end{clearBox}}    


%----------------------------------------------------------------------------------------
%	REMARK ENVIRONMENT
%----------------------------------------------------------------------------------------

\newenvironment{remark}{\par\vskip10pt\small % Vertical white space above the remark and smaller font size
\begin{list}{}{
\leftmargin=35pt % Indentation on the left
\rightmargin=25pt}\item\ignorespaces % Indentation on the right
\makebox[-2.5pt]{\begin{tikzpicture}[overlay]
\node[draw=color!60,line width=1pt,circle,fill=color!25,font=\sffamily\bfseries,inner sep=2pt,outer sep=0pt] at (-15pt,0pt){\textcolor{color}{N}};\end{tikzpicture}} % Orange R in a circle
\advance\baselineskip -1pt}{\end{list}\vskip5pt} % Tighter line spacing and white space after remark

%----------------------------------------------------------------------------------------
%	SECTION NUMBERING IN THE MARGIN
%----------------------------------------------------------------------------------------

\makeatletter
\renewcommand{\@seccntformat}[1]{\llap{\textcolor{color}{\csname the#1\endcsname}\hspace{1em}}}                    
\renewcommand{\section}{\@startsection{section}{1}{\z@}
{-4ex \@plus -1ex \@minus -.4ex}
{1ex \@plus.2ex }
{\normalfont\large\sffamily\bfseries}}
\renewcommand{\subsection}{\@startsection {subsection}{2}{\z@}
{-3ex \@plus -0.1ex \@minus -.4ex}
{0.5ex \@plus.2ex }
{\normalfont\sffamily\bfseries}}
\renewcommand{\subsubsection}{\@startsection {subsubsection}{3}{\z@}
{-2ex \@plus -0.1ex \@minus -.2ex}
{0.2ex \@plus.2ex }
{\normalfont\small\sffamily\bfseries}}                        
\renewcommand\paragraph{\@startsection{paragraph}{4}{\z@}
{-2ex \@plus-.2ex \@minus .2ex}
{0.1ex}
{\normalfont\small\sffamily\bfseries}}

%----------------------------------------------------------------------------------------
%	CHAPTER HEADINGS
%----------------------------------------------------------------------------------------

\newcommand{\thechapterimage}{}
\newcommand{\chapterimage}[1]{\renewcommand{\thechapterimage}{#1}}
\def\thechapter{\arabic{chapter}}
\def\@makechapterhead#1{
\thispagestyle{empty}
{\centering \normalfont\sffamily
\ifnum \c@secnumdepth >\m@ne
\if@mainmatter
\startcontents
\begin{tikzpicture}[remember picture,overlay]
\node at (current page.north west)
{\begin{tikzpicture}[remember picture,overlay]

\node[anchor=north west,inner sep=0pt] at (0,0) {\includegraphics[width=\paperwidth]{\thechapterimage}};

%Commenting the 3 lines below removes the small contents box in the chapter heading
\draw[fill=white,opacity=1] (1cm, 0) rectangle (8cm,-7cm);
\node[anchor=north west] at (1cm,.25cm) {\parbox[t][8cm][t]{6.5cm}{\huge\bfseries\flushleft \printcontents{l}{1}{\setcounter{tocdepth}{2}}}};
%\node[anchor=north west] at (1cm,.25cm) {\parbox[t][8cm][t]{6.5cm}{\huge\bfseries\flushleft \printcontents{l}{1}{\setcounter{tocdepth}{1}}}};
\draw[anchor=west] (5cm,-9cm) node [rounded corners=25pt,fill=white,fill opacity=1,text opacity=1,draw=color,draw opacity=1,line width=2pt,inner sep=15pt]{\huge\sffamily\bfseries\textcolor{black}{\thechapter\ ---\ #1\vphantom{plPQq}\makebox[22cm]{}}};

\end{tikzpicture}};
\end{tikzpicture}}\par\vspace*{230\p@}
\fi
\fi
}
\def\@makeschapterhead#1{
\thispagestyle{empty}
{\centering \normalfont\sffamily
\ifnum \c@secnumdepth >\m@ne
\if@mainmatter
%\startcontents
\begin{tikzpicture}[remember picture,overlay]
\node at (current page.north west)
{\begin{tikzpicture}[remember picture,overlay]
%\node[anchor=north west,inner sep=0pt] at (0,0) {\includegraphics[width=\paperwidth]{\thechapterimage}};
%Commenting the 3 lines below removes the small contents box in the chapter heading
%\draw[fill=white,opacity=1] (1cm,0) rectangle (8cm,-7cm);
%\node[anchor=north west] at (1cm,.25cm) {\parbox[t][8cm][t]{6.5cm}{\huge\bfseries\flushleft \printcontents{l}{1}{\setcounter{tocdepth}{2}}}};
\draw[anchor=west] (-0.2cm,-2.5cm) node [fill=white,fill opacity=1,text opacity=1,draw=color,draw opacity=1,line width=2pt,inner sep=15pt, inner xsep=60pt]{\huge\sffamily\bfseries\textcolor{black}{  \ #1\vphantom{plPQq}\makebox[22cm]{}}};
\end{tikzpicture}};
\end{tikzpicture}}\par\vspace*{50\p@}
\fi
\fi
}
\makeatother % Insert the commands.tex file which contains the majority of the structure behind the template

%----------------------------------------------------------------------------------------

\begin{document}
\graphicspath{ {../../Images/}}
\fi


\chapterimage{Headers/chapter_header.pdf}
\chapter[Additional elements from the Legrand orange Book]{Additional elements
\label{chap:legrand}}
%\begin{refsection}
%------------------------------------------------
%\listoftodos


\section{Theorems}\index{Theorems}

This is an example of theorems.

\subsection{Several equations}\index{Theorems!Several Equations}
This is a theorem consisting of several equations.

\begin{theorem}[Name of the theorem]
In $E=\mathbb{R}^n$ all norms are equivalent. It has the properties:
\begin{align}
& \big| ||\mathbf{x}|| - ||\mathbf{y}|| \big|\leq || \mathbf{x}- \mathbf{y}||\\
&  ||\sum_{i=1}^n\mathbf{x}_i||\leq \sum_{i=1}^n||\mathbf{x}_i||\quad\text{where $n$ is a finite integer}
\end{align}
\end{theorem}

\subsection{Single Line}\index{Theorems!Single Line}
This is a theorem consisting of just one line.

\begin{theorem}
A set $\mathcal{D}(G)$ in dense in $L^2(G)$, $|\cdot|_0$. 
\end{theorem}

%------------------------------------------------

\section{Definitions}\index{Definitions}

This is an example of a definition. A definition could be mathematical or it could define a concept.

\begin{definition}[Definition name]
Given a vector space $E$, a norm on $E$ is an application, denoted $||\cdot||$, $E$ in $\mathbb{R}^+=[0,+\infty[$ such that:
\begin{align}
& ||\mathbf{x}||=0\ \Rightarrow\ \mathbf{x}=\mathbf{0}\\
& ||\lambda \mathbf{x}||=|\lambda|\cdot ||\mathbf{x}||\\
& ||\mathbf{x}+\mathbf{y}||\leq ||\mathbf{x}||+||\mathbf{y}||
\end{align}
\end{definition}

%------------------------------------------------

\section{Notations}\index{Notations}

\begin{notation}
Given an open subset $G$ of $\mathbb{R}^n$, the set of functions $\varphi$ are:
\begin{enumerate}
\item Bounded support $G$;
\item Infinitely differentiable;
\end{enumerate}
a vector space is denoted by $\mathcal{D}(G)$. 
\end{notation}

%------------------------------------------------

\section{Remarks}\index{Remarks}

This is an example of a remark.

\begin{remark}
The concepts presented here are now in conventional employment in mathematics. Vector spaces are taken over the field $\mathbb{K}=\mathbb{R}$, however, established properties are easily extended to $\mathbb{K}=\mathbb{C}$.
\end{remark}

%------------------------------------------------

\section{Corollaries}\index{Corollaries}

This is an example of a corollary.

\begin{corollary}[Corollary name]
The concepts presented here are now in conventional employment in mathematics. Vector spaces are taken over the field $\mathbb{K}=\mathbb{R}$, however, established properties are easily extended to $\mathbb{K}=\mathbb{C}$.
\end{corollary}

%------------------------------------------------

\section{Propositions}\index{Propositions}

This is an example of propositions.

\subsection{Several equations}\index{Propositions!Several Equations}

\begin{proposition}[Proposition name]
It has the properties:
\begin{align}
& \big| ||\mathbf{x}|| - ||\mathbf{y}|| \big|\leq || \mathbf{x}- \mathbf{y}||\\
&  ||\sum_{i=1}^n\mathbf{x}_i||\leq \sum_{i=1}^n||\mathbf{x}_i||\quad\text{where $n$ is a finite integer}
\end{align}
\end{proposition}

\subsection{Single Line}\index{Propositions!Single Line}

\begin{proposition} 
Let $f,g\in L^2(G)$; if $\forall \varphi\in\mathcal{D}(G)$, $(f,\varphi)_0=(g,\varphi)_0$ then $f = g$. 
\end{proposition}

%------------------------------------------------

\section{Examples}\index{Examples}

This is an example of examples.

\subsection{Equation and Text}\index{Examples!Equation and Text}

\begin{example}
Let $G=\{x\in\mathbb{R}^2:|x|<3\}$ and denoted by: $x^0=(1,1)$; consider the function:
\begin{equation}
f(x)=\left\{\begin{aligned} & \mathrm{e}^{|x|} & & \text{si $|x-x^0|\leq 1/2$}\\
& 0 & & \text{si $|x-x^0|> 1/2$}\end{aligned}\right.
\end{equation}
The function $f$ has bounded support, we can take $A=\{x\in\mathbb{R}^2:|x-x^0|\leq 1/2+\epsilon\}$ for all $\epsilon\in\intoo{0}{5/2-\sqrt{2}}$.
\end{example}

\subsection{Paragraph of Text}\index{Examples!Paragraph of Text}

\begin{example}[Example name]
\lipsum[2]
\end{example}

%------------------------------------------------

\section{Exercises}\index{Exercises}

This is an example of an exercise.

\begin{exercise}
This is a good place to ask a question to test learning progress or further cement ideas into students' minds.
\end{exercise}

%------------------------------------------------

\section{Problems}\index{Problems}

\begin{problem}
What is the average airspeed velocity of an unladen swallow?
\end{problem}

%------------------------------------------------

\section{Vocabulary}\index{Vocabulary}

Define a word to improve a students' vocabulary.

\begin{vocabulary}[Word]
Definition of word.
\end{vocabulary}



%\fi
%------------------------------------------------
%\printbibliography[heading=subbibliography]
%\end{refsection}

\iffalse
\begingroup
\let\enotesize\normalsize
\renewcommand\notesname{Notas finais}
\renewcommand\enoteformat{\noindent[\theenmark]\hspace{0.06in}}
\label{notas}
\theendnotes
\endgroup

\end{document}
\fi




%\chapterimage{./Headers/general} % Chapter heading image
%\chapter{Apêndices}

\chapter*{Notas finais}
\begingroup
\let\enotesize\normalsize
\renewcommand\notesname{}
\renewcommand\enoteformat{\noindent[\theenmark]\hspace{0.06in}}
\label{notas}
\theendnotes
\endgroup

\begin{appendix}
 
\end{appendix}



%----------------------------------------------------------------------------------------
%	BIBLIOGRAPHY
%----------------------------------------------------------------------------------------

%\chapter*{Bibliography}
%\addcontentsline{toc}{chapter}{\textcolor{color}{Bibliography}}
%\section*{Books}
%\addcontentsline{toc}{section}{Books}
%\printbibliography[heading=bibempty,type=book]
%\section*{Articles}
%\addcontentsline{toc}{section}{Articles}
%\printbibliography[heading=bibempty,type=article]
%\printbibliography

%----------------------------------------------------------------------------------------
%	INDEX
%----------------------------------------------------------------------------------------
%\chapter*{Índice remissivo}
%\cleardoublepage
%\setlength{\columnsep}{0.75cm}
%\addcontentsline{toc}{chapter}{\textcolor{color}{Índice remissivo}}
%\printindex

%----------------------------------------------------------------------------------------

\end{document}